%%%%%%%%%%%%%%%%%%%%%%%%%%%%%%%%%%%%%%%%%%%%%%%%%%%%%%%%%%%%%%%%%%%%%%
% Lớp 11 | Đạo hàm | Đạo hàm
% Dạng 1D7?1-1: Tính đạo hàm bằng định nghĩa
% Số câu: 71  (sinh tự động bởi tach_cau.py; từ nay thầy sửa trực tiếp file này)
%%%%%%%%%%%%%%%%%%%%%%%%%%%%%%%%%%%%%%%%%%%%%%%%%%%%%%%%%%%%%%%%%%%%%%

%%%%%%%%%%%%%%%% TRẮC NGHIỆM 4 PHƯƠNG ÁN %%%%%%%%%%%%%%%%

% nguồn: Ok_Toan11_TracNghiem_Lan2
\begin{ex}%[1D7N1-1]%[ma:4ea7afd2]
	Cho hàm số $ y=f(x)  $ có đạo hàm tại điểm $ x_{0} $. Tìm khẳng định đúng trong các khẳng định sau
	\choice 
		{\True $  f'\left(x_{0}\right)=\lim\limits_{x \rightarrow x_{0}} \dfrac{f(x)-f\left(x_{0}\right)}{x-x_{0}} $}
		 {$ f'\left(x_{0}\right)=\lim\limits_{x \rightarrow x_{0}} \dfrac{f(x)+f\left(x_{0}\right)}{x-x_{0}} $}
		{$ f'\left(x_{0}\right)=\lim\limits_{x \rightarrow x_{0}}  \dfrac{f(x)-f\left(x_{0}\right)}{x+x_{0}} $}
		{$ f'\left(x_{0}\right)=\lim\limits_{x \rightarrow x_{0}}  \dfrac{f(x)+f\left(x_{0}\right)}{x+x_{0}} $}
	 \loigiai{
		Theo định nghĩa đạo hàm ta có $ f'\left(x_{0}\right)=\lim\limits_{x \rightarrow x_{0}} \dfrac{f(x)-f\left(x_{0}\right)}{x-x_{0}}  $.
	 }
\end{ex}

% nguồn: Ok_Toan11_TracNghiem_Lan2
\begin{ex}%[1D7N1-1]%[ma:95b927aa]
	Cho hàm số $ y=f(x) $ có đạo  hàm thỏa mãn $ f'(6)=2 $. Giá trị của biểu thức $ \lim\limits_{x \rightarrow 6} \dfrac{f(x)-f(6)}{x-6} $ bằng
	\choice 
		 {12}
		{\True 2} 
		{$\dfrac{1}{3}$} 
		 {$\dfrac{1}{2}$}
	\loigiai{ 
		Hàm số $ y=f(x) $ có tập xác định là $ \mathscr{D} $ và $ x_{0} \in \mathscr{D} $.\\ Nếu tồn tại giới hạn $ \lim\limits_{x \rightarrow x_{0}} \dfrac{f(x)-f\left(x_{0}\right)}{x-x_{0}} $ thì giới hạn gọi là  đạo hàm của hàm số tại $ x_{0} $.\\
		Vậy kết quả của  biểu thức $ \lim\limits_{x \rightarrow 6} \dfrac{f(x)-f(6)}{x-6}=f'(6)=2 $.
	} 
\end{ex}

% nguồn: Ok_Toan11_TracNghiem_Lan2
\begin{ex}%[1D7N1-1]%[ma:95ff4277]
	Cho hàm số $  y=f(x) $ xác định trên $ \mathbb{R} $ thỏa mãn $ \lim\limits_{x \rightarrow 3} \dfrac{f(x)-f(3)}{x-3}=2 $. Kết quả đúng là
	 \choice
		{$ f'(2)=3  $}
		{$ f'(x)=2  $}
		{$ f'(x)=3  $}
		{\True $ f'(3)=2 $} 
	 \loigiai{
		Theo định  nghĩa đạo hàm của hàm số tại một điểm ta có
		 \[
		\lim\limits_{x \rightarrow 3} \dfrac{f(x)-f(3)}{x-3}=2=f'(3). 
		 \]
	} 
\end{ex}

% nguồn: Ok_TracNghiemLop11_Lan1
\begin{ex}%[1D7N1-1]%[ma:dc0fecc5]
	Cho hàm  số $y=f(x)$ xác định trên $\mathbb{R}$ thỏa mãn $\lim\limits_{x\to 4}\dfrac{f(x)-f(4)}{x-4}=5$. Kết quả đúng là
	\choice 
		{$f'(2)=3$} 
		 {$f'(x)=2$}
		 {$f'(x)=5$}
		 {\True $f'(4)=5$}
	\loigiai{ 
		Theo định nghĩa đạo hàm của hàm  số tại một điểm ta có
	$\lim\limits_{x\to 4}\dfrac{f(x)-f(4)}{x-4}=5=f'(4)$.} 
\end{ex}

% nguồn: Ok_TracNghiemLop11_Lan1
\begin{ex}%[1D7N1-1]%[ma:fd430cbd]
	Phát biểu nào trong các phát biểu sau  là đúng?
	\choice 
		{Nếu hàm số $y=f(x)$ có đạo hàm trái tại $x_{0}$ thì nó liên tục  tại điểm đó}
		{Nếu hàm số $y=f(x)$ có đạo hàm phải tại $x_{0}$ thì nó liên tục  tại điểm đó}
		{Nếu hàm  số $y=f(x)$ có đạo hàm tại $x_{0}$ thì nó liên tục tại điểm $-x_{0}$}
		{\True Nếu  hàm số $y=f(x)$ có đạo hàm tại $x_{0}$ thì nó liên tục tại điểm đó}
	\loigiai{ 
		Ta có định  lí sau\\
		Nếu hàm  số $y=f(x)$ có đạo hàm tại $x_{0}$ thì nó liên tục tại điểm đó.
	 }
\end{ex}

% nguồn: Ok_TracNghiemLop11_Lan1
\begin{ex}%[1D7N1-1]%[ma:62b2a98f]
	Cho hàm  số $y=f(x)$ xác định trên $\mathbb{R}$ thỏa mãn $\displaystyle\underset{x\to 3}{\mathop{\lim}}\dfrac{f(x)-f\left(3\right)}{x-3}=2$. Kết quả đúng là
	 \choice
		{$f'\left(2\right)=3$} 
		{$f'\left(x\right)=2$} 
		{$f'\left(x\right)=3$} 
		{\True  $f'\left(3\right)=2$}
	\loigiai{ 
		Theo  định nghĩa đạo hàm của hàm số tại một điểm ta có
	$$\underset{x\to  3}{\mathop{\lim}}\dfrac{f(x)-f\left(3\right)}{x-3}=2=f'\left(3\right).$$}
\end{ex}

% nguồn: Ok_TracNghiemLop11_Lan1
\begin{ex}%[1D7N1-1]%[ma:1d9ba4a7]
	Cho hàm số $y=f(x)$ có đạo hàm tại điểm $x_0$. Tìm khẳng định đúng trong các  khẳng định sau
	\choice 
		{\True  $f'(x_0)=\lim\limits_{x\to x_0}\dfrac{f(x)-f(x_0)}{x-x_0}$}
		{$f'(x_0)=\lim\limits_{x\to x_0}\dfrac{f(x)+f(x_0)}{x-x_0}$} 
		 {$f'(x_0)=\lim\limits_{x\to x_0}\dfrac{f(x)-f(x_0)}{x+x_0}$}
		{$f'(x_0)=\lim\limits_{x\to  x_0}\dfrac{f(x)+f(x_0)}{x+x_0}$}
	\loigiai{ 
	Theo  định nghĩa đạo hàm ta có $f'(x_0)=\lim\limits_{x\to x_0}\dfrac{f(x)-f(x_0)}{x-x_0}$.}
\end{ex}

% nguồn: Ok_TracNghiemLop11_Lan1
\begin{ex}%[1D7N1-1]%[ma:2cd8cc4d]
	Cho hàm số $ y=f(x)$  có đạo hàm thỏa mãn $f'(6)=2$. Giá trị của biểu thức $\lim\limits_{x\to 6}\dfrac{f(x)-f(6)}{x-6}$ bằng
	\choice 
		 {$12$}
		{\True $2$} 
		 {$\dfrac{1}{3}$}
		{$\dfrac{1}{2}$} 
	 \loigiai{
		Ta  có $\lim\limits_{x\to 6}\dfrac{f(x)-f(6)}{x-6}=f'(6)=2$.
	} 
\end{ex}

% nguồn: Ok_Toan11_TracNghiem_Lan2
\begin{ex}%[1D7H1-1]%[ma:e57b5874]
	Cho hàm số $ f(x)=\heva{&\dfrac{\sqrt{3 x+1}-2 x}{x-1} \text { khi } x \neq 1 \\& \dfrac{-5}{4} \text  { khi } x=1}$. Khi đó $ f'(1) $
	 \choice
		{Không tồn  tại}
		 {$0$}
		 {\True $\dfrac{3}{8} $}
		{$  -\dfrac{9}{64} $}
	\loigiai{ 
		 Ta có
		\begin{eqnarray*}
			\lim\limits_{x \rightarrow 1} f(x)&=&\lim\limits_{x  \rightarrow 1} \dfrac{\sqrt{3 x+1}-2 x}{x-1}\\
			&=&\lim\limits_{x \rightarrow 1} \dfrac{3 x+1-4 x^{2}}{(x-1)(\sqrt{3 x+1}+2 x)}\\ 
			&=&\lim\limits_{x  \rightarrow 1} \dfrac{(1-x)(4x+1)}{(x-1)(\sqrt{3 x+1}+2 x)}=\dfrac{-5}{4}=f(1).
		\end{eqnarray*}
		Hàm số  liên tục tại $ x=1 $. Ta có
		 \[
		\begin{aligned}
			f'(1) & =\lim\limits_{x  \rightarrow 1} \dfrac{f(x)-f(1)}{x-1}=\lim\limits_{x \rightarrow 1} \dfrac{\dfrac{\sqrt{3 x+1}-2 x}{x-1}+\dfrac{5}{4}}{x-1}=\lim\limits_{x \rightarrow 1} \dfrac{4 \sqrt{3 x+1}-8x+5}{4(x-1)^{2}} \\
			& =\lim\limits_{x \rightarrow 1} \dfrac{4 \sqrt{3 x+1}-8x+5}{4(x-1)^{2}} = \lim\limits_{x \rightarrow 1} \dfrac{12(x-1) -  8(x-1)^2}{8(x-1)(\sqrt{3x+1} +2x)}\\
			& = \lim\limits_{x \rightarrow 1} \dfrac{12  - 8(x-1)}{8(\sqrt{3x+1} +2x)} = \dfrac{12}{32} = \dfrac{3}{8}.
		\end{aligned}
		 \]
	 }
\end{ex}

% nguồn: Ok_Toan11_TracNghiem_Lan2
\begin{ex}%[1D7H1-1]%[ma:0cef8bc9]
	Cho hàm số  $ f(x) $ xác định bởi $f(x)=\heva{&\dfrac{\sqrt{4 x^{2}+1}-1}{x} \text { khi } x \neq 0 \\& 0 \text { khi } x=0}$. Giá trị $ f'(0) $ bằng
	\choice 
		 {\True $2$}
		 {$0$}
		 {$\dfrac{1}{2}$}
		{Không  tồn tại}
	 \loigiai{
		 Tập xác định $ \mathscr{D}=\mathbb{R} $.\\
		Ta có 
		\begin{eqnarray*}
			 \lim\limits_{x \rightarrow 0} \dfrac{f(x)-f(0)}{x-0}&=&\lim\limits_{x \rightarrow 0} \dfrac{\sqrt{4 x^{2}+1}-1}{x^{2}}\\
			&=&\lim\limits_{x \rightarrow 0} \dfrac{4  x^{2}}{x^{2}\left(\sqrt{4 x^{2}+1}+1\right)}\\
			&=&\lim\limits_{x \rightarrow  0} \dfrac{4}{\sqrt{4 x^{2}+1}+1}=2.
		\end{eqnarray*}
		Vậy $ f'(0)=2 $. 
	 }
\end{ex}

% nguồn: Ok_Toan11_TracNghiem_Lan2
\begin{ex}%[1D7H1-1]%[ma:c3e25676]
	Cho hàm số $ f(x)=\heva{&(x-1)^{2} \text { khi } x \geq 0 \\& -x^{2} \text { khi } x<0}$ có đạo hàm tại điểm $ x_{0}=0 $ là 
	 \choice
		 {$f'(0)=0$}
		{$f'(0)=1$} 
		{$f'(0)=-2$} 
		 {\True Không tồn tại}
	\loigiai{ 
		Ta có 
		\begin{itemize}
			\item $f(0)=1$. 
			\item $\lim\limits_{x  \rightarrow 0^{+}} f(x)=\lim\limits_{x \rightarrow 0^{+}}(x-1)^{2}=1$.
			\item $\lim\limits_{x \rightarrow 0^{-}} f(x)=\lim\limits_{x \rightarrow 0^{-}}\left(-x^{2}\right)=0$. 
		\end{itemize}
		Ta thấy $ f(0)=\lim\limits_{x \rightarrow 0^{+}} f(x) \neq \lim\limits_{x \rightarrow 0^{-}} f(x) $ nên hàm số không liên tục tại  $ x_{0}=0 $.\\
		Vậy hàm số không có đạo  hàm tại $ x_{0}=0 $.
	} 
\end{ex}

% nguồn: Ok_Toan11_TracNghiem_Lan2
\begin{ex}%[1D7H1-1]%[ma:203f6f12]
	Cho hàm số $ f(x)=\heva{&\dfrac{3-\sqrt{4-x}}{4}  \text { khi } x \neq 0 \\& \dfrac{1}{4} \text { khi } x=0}$. Khi đó $ f'(0) $ bằng
	 \choice
		{$\dfrac{1}{4}$} 
		{\True $\dfrac{1}{16}$} 
		{$\dfrac{1}{32}$} 
		{Không tồn  tại}
	 \loigiai{
		Với $ x \neq 0  $ xét
		\begin{eqnarray*}
			\lim\limits_{x \rightarrow 0} \dfrac{f(x)-f(0)}{x-0}&=&\lim\limits_{x \rightarrow  0} \dfrac{\dfrac{3-\sqrt{4-x}}{4}-\dfrac{1}{4}}{x}\\
			&=&\lim\limits_{x \rightarrow 0} \dfrac{2-\sqrt{4-x}}{4 x}=\lim\limits_{x \rightarrow 0} \dfrac{4-(4-x)}{4 x(2+\sqrt{4-x})} \\ 
			&=&\lim\limits_{x \rightarrow 0}  \dfrac{1}{4(2+\sqrt{4-x})}=\dfrac{1}{4(2+\sqrt{4-0})}=\dfrac{1}{16}.
		\end{eqnarray*}
		 Suy ra $f'(0)=\dfrac{1}{16}$.
	 }
\end{ex}

% nguồn: Ok_Toan11_TracNghiem_Lan2
\begin{ex}%[1D7H1-1]%[ma:b0eee598]
	Cho hàm số $ y=f(x)=\heva{&x^{2}+1\text{ khi }  x \geq 1 \\& 2 x \text{ khi } x<1}$. Mệnh đề \textbf{sai} là
	\choice 
		{$  f'(1)=2 $}
		{\True $ f $ không có đạo hàm tại $ x_{0}=1  $}
		{$  f'(0)=2 $}
		{$ f'(2)=4 $} 
	\loigiai{ 
		Ta  có
		\begin{itemize}
			\item $  \lim\limits_{x \rightarrow 1^{-}} \dfrac{f(x)-f(1)}{x-1}=\lim\limits_{x \rightarrow 1^{-}} \dfrac{2 x-2}{x-1}=2 $.
			\item $\lim\limits_{x \rightarrow 1^{+}} \dfrac{f(x)-f(1)}{x-1}=\lim\limits_{x \rightarrow 1^{+}} \dfrac{x^{2}+1-2}{x-1}=\lim\limits_{x \rightarrow  1^{+}}(x+1)=2$.
		\end{itemize}
		Suy  ra $ f'\left(1^{-}\right)=f'\left(1^{+}\right)=f'(1)=2 $.\\
		Vậy hàm số có đạo hàm  tại $ x_{0}=1 $.
	} 
\end{ex}

% nguồn: Ok_Toan11_TracNghiem_Lan2
\begin{ex}%[1D7H1-1]%[ma:2aa1aa3e]
	Cho hàm số $ y=f(x) $ có đạo hàm tại điểm $ x_{0}=2  $. Khi đó $ \lim\limits_{x \rightarrow 2} \dfrac{2 f(x)-x f(2)}{x-2} $ bằng
	\choice 
		 {$0$}
		{$f'(2)$} 
		{\True  $2 f'(2)-f(2)$}
		 {$f(2)-2 f'(2)$}
	 \loigiai{
		Do hàm số $ y=f(x) $ có đạo hàm tại điểm $ x_{0}=2 $ suy ra  $ \lim\limits_{x \rightarrow 2} \dfrac{f(x)-f(2)}{x-2}=f'(2) $.\\
		 Ta có
		\begin{eqnarray*}
			 I &= &\lim\limits_{x \rightarrow 2} \dfrac{2 f(x)-x f(2)}{x-2} \\
			&=&\lim\limits_{x \rightarrow 2} \dfrac{2 f(x)-2 f(2)+2  f(2)-x f(2)}{x-2} \\
			&=&\lim\limits_{x \rightarrow  2} \dfrac{2(f(x)-f(2))}{x-2}-\lim\limits_{x \rightarrow 2} \dfrac{f(2)(x-2)}{x-2} \\
			&=&2 f'(2)-f(2). 
		\end{eqnarray*}
	 }
\end{ex}

% nguồn: Ok_Toan11_TracNghiem_Lan2
\begin{ex}%[1D7H1-1]%[ma:dfb029d9]
	Cho hàm số $ y=\heva{&x^{2}+a x+b \text { khi } x \geq 2 \\& x^{3}-x^{2}-8  x+10 \text { khi } x<2}$. Biết hàm số có đạo hàm tại điểm $ x=2 $. Giá trị của $ a^{2}+b^{2} $ bằng
	\choice 
		{\True  $20$}
		{$17$} 
		{$18$} 
		{$25$} 
	\loigiai{ 
		Ta có $y'=\heva{&2 x+a \text { khi } x \geq 2 \\& 
		3  x^{2}-2 x-8 \text { khi } x<2.}$\\
		Hàm số có đạo hàm tại  điểm $ x=2 \Leftrightarrow 4+a=0 \Leftrightarrow a=-4 $.\\
		Mặt khác hàm số có  đạo hàm tại điểm $ x=2 $ thì hàm số liên tục tại điểm $ x=2 $.\\
		Suy ra $$ \lim\limits_{x \rightarrow 2^{+}} f(x)=\lim\limits_{x \rightarrow 2^{-}} f(x)=f(2)\Leftrightarrow  4+2 a+b=-2 \Leftrightarrow b=2.$$
		 Vậy $ a^{2}+b^{2}=20 $.
	} 
\end{ex}

% nguồn: Ok_Toan11_TracNghiem_Lan2
\begin{ex}%[1D7H1-1]%[ma:c4be8a27]
	Cho hàm số $ f(x)=\heva{&a x^{2}+b x \text { khi } x \geq 1 \\& 2 x-1 \text { khi }  x<1}$. Để hàm số đã cho có đạo hàm tại $ x=1 $ thì $ 2 a+b $ bằng
	\choice 
		 {\True $2$}
		{$5$} 
		{$-2$} 
		{$-5$} 
	 \loigiai{
		Ta có 
		\begin{itemize}
			\item $\lim\limits_{x \rightarrow 1^{-}} \dfrac{f(x)-f(1)}{x-1}=\lim\limits_{x \rightarrow 1^{-}} \dfrac{2 x-1-1}{x-1}=2$. 
			\item 
		\end{itemize}
		\begin{eqnarray*}
			\lim\limits_{x \rightarrow 1^{+}} \dfrac{f(x)-f(1)}{x-1}&=&\lim\limits_{x  \rightarrow 1^{+}} \dfrac{a x^{2}+b x-a-b}{x-1}=\lim\limits_{x \rightarrow 1^{+}} \dfrac{a\left(x^{2}-1\right)+b(x-1)}{x-1}\\
			&=&\lim\limits_{x \rightarrow 1^{+}} \dfrac{(x-1)[a(x+1)+b]}{x-1}\\ 
			&=&=\lim\limits_{x \rightarrow 1^{+}}[a(x+1)+b]=2  a+b.
		\end{eqnarray*}
		Theo yêu cầu bài  toán $ \lim\limits_{x \rightarrow 1^{-}} \dfrac{f(x)-f(1)}{x-1}=\lim\limits_{x \rightarrow 1^{+}} \dfrac{f(x)-f(1)}{x-1} \Leftrightarrow 2 a+b=2 $.
	} 
\end{ex}

% nguồn: Ok_Toan11_TracNghiem_Lan2
\begin{ex}%[1D7H1-1]%[ma:f84e92ca]
	Cho hàm số $ f(x)=\dfrac{3 x}{1+|x|} $. Tính $ f'(0) $. 
	 \choice
		 {$ f'(0)=0 $}
		{$ f'(0)=1 $} 
		 {$ f'(0)=\dfrac{1}{3} $}
		{\True $ f'(0)=3 $} 
	\loigiai{ 
		Ta  có $ f'(0)=\lim\limits_{x \rightarrow 0} \dfrac{f(x)-f(0)}{x}=\lim\limits_{x \rightarrow 0} \dfrac{3}{1+|x|} $.\\
		Mà 
		\begin{itemize}
			\item $\lim\limits_{x \rightarrow 0^{+}} \dfrac{3}{1+|x|}=\lim\limits_{x \rightarrow  0^{+}} \dfrac{3}{1+x}=3$.
			\item  $\lim\limits_{x \rightarrow 0^{-}} \dfrac{3}{1+|x|}=\lim\limits_{x \rightarrow 0^{-}} \dfrac{3}{1-x}=3$.
		\end{itemize}
		Suy ra $\lim\limits_{x \rightarrow 0^{+}} \dfrac{3}{1+|x|}=\lim\limits_{x \rightarrow 0^{-}} \dfrac{3}{1+|x|}=3  $.\\
		 Do đó $f'(0)=\lim\limits_{x \rightarrow 0} \dfrac{3}{1+|x|}=3 $.\\
		 Vậy $ f'(0)=3 $.
	 }
\end{ex}

% nguồn: Ok_Toan11_TracNghiem_Lan2
\begin{ex}%[1D7H1-1]%[ma:a69225e5]
	Tính đạo hàm của hàm số số $ y=x(x-1)(x-2) \ldots(x-2021) $ tại điểm $  x=0 $.
	 \choice
		{$f'(0)=0$} 
		 {$f'(0)=2021!$}
		 {$f'(0)=2021$}
		 {\True $f'(0)=-2021!$}
	\loigiai{ 
		Ta có 
		\begin{eqnarray*}
			f'(0)&=&\lim\limits_{x \rightarrow 0} \dfrac{f(x)-f(0)}{x-0}=\lim\limits_{x  \rightarrow 0} \dfrac{x(x-1)(x-2) \ldots(x-2021)}{x} \\
			&=&\lim\limits_{x \rightarrow 0}(x-1)(x-2) \ldots(x-2021)\\ 
			&=&(-1) \cdot(-2)  \ldots(-2021)=-2021!.
		\end{eqnarray*}
	 }
\end{ex}

% nguồn: Ok_TracNghiemLop11_Lan1
\begin{ex}%[1D7H1-1]%[ma:032d9802]
	] 
	Tính đạo hàm của hàm  số
	$$ 
	f(x)=\heva{&\dfrac{\sqrt{x^3-2 x^2+x+1}-1}{x-1}, &&x  \neq 1 \\ &0, &&x=1}
	$$  tại điểm $x _0=1$.
	\choice 
		{$ \dfrac{1}{3}  $}
		{$ \dfrac{1}{5}  $}
		{\True $ \dfrac{1}{2}  $}
		{$ \dfrac{1}{4}  $}
	\loigiai{ 
		Ta có $f(1)=0$ và 
		\begin{align*}
			\lim\limits_{x  \rightarrow 1} \dfrac{f(x)-f(1)}{x-1}
			&=\lim\limits_{x \rightarrow  1} \dfrac{\sqrt{x^3-2 x^2+x+1}-1}{(x-1)^2} \\
			& =\lim\limits_{x \rightarrow 1} \dfrac{x^3-2 x^2+x}{(x-1)^2 \cdot\left(\sqrt{x^3-2  x^2+x+1}+1\right)} \\
			& =\lim\limits_{x \rightarrow 1} \dfrac{x(x-1)^2}{(x-1)^2 \cdot\left(\sqrt{x^3-2  x^2+x+1}+1\right)} \\
			& =\lim\limits_{x \rightarrow 1} \dfrac{x}{\sqrt{x^3-2 x^2+x+1}+1} \\ 
			&=\dfrac{1}{2}. 
		\end{align*}
		 Vậy $f'(1)=\dfrac{1}{2}$.
	} 
\end{ex}

% nguồn: Ok_TracNghiemLop11_Lan1
\begin{ex}%[1D7H1-1]%[ma:6ef74108]
	Cho hàm số $y=f(x)$ có đạo hàm thỏa  mãn $f'(6)=2$. Giá trị của biểu thức $\lim\limits_{x \rightarrow 6} \dfrac{f(x)-f(6)}{x-6}$ bằng
	\choice 
		 {$ 12 $}
		{\True $ 2  $}
		 {$\dfrac{1}{3}$}
		 {$\dfrac{1}{2}$}
	\loigiai{ 
		Hàm số $y=f(x)$ có tập xác định là $\mathscr{D}$ và $x_0 \in \mathscr{D}$. 
		 \\
		Nếu tồn tại giới hạn $\lim\limits_{x \rightarrow x_0} \dfrac{f(x)-f\left(x_0\right)}{x-x_0}$ thì giới hạn gọi là  đạo hàm của hàm số tại $x_0$.
		 \\
		Vậy  kết quả của biểu thức
		 $$
		\lim\limits_{x \rightarrow 6}  \dfrac{f(x)-f(6)}{x-6}=f'(6)=2.
		 $$
	 }
\end{ex}

% nguồn: Ok_TracNghiemLop11_Lan1
\begin{ex}%[1D7H1-1]%[ma:a990f857]
	Tính đạo  hàm của hàm số $y=\dfrac{2 x-1}{x+1}$ tại $x =3$
	\choice 
		{$ \dfrac{1}{6} $} 
		{\True $  \dfrac{3}{16} $}
		 {$ \dfrac{2}{9} $}
		{$ \dfrac{4}{5} $} 
	\loigiai{ 
		 $$
		 \lim\limits_{x \rightarrow 3} \dfrac{f(x)-f(3)}{x-3}
		=\lim\limits_{x \rightarrow 3} \dfrac{\dfrac{2 x-1}{x+1}-\dfrac{5}{4}}{x-3} 
		=\lim\limits_{x \rightarrow 3}  \dfrac{3(x-3)}{(x-3)(x+1) 4}
		=\lim\limits_{x  \rightarrow 3} \dfrac{3}{(x+1) 4}
		=\dfrac{3}{16}. 
		 $$
		Vậy hàm số có đạo hàm tại  $x=3$ là $f'(3)=\dfrac{3}{16}$.
	} 
\end{ex}

% nguồn: Ok_TracNghiemLop11_Lan1
\begin{ex}%[1D7H1-1]%[ma:3dbd6792]
	Cho hàm số $y=f(x)$ có đạo hàm tại điểm $x_0$. Tìm khẳng định đúng trong các  khẳng định sau
	 \choice
		{\True  $f'\left(x_0\right)=\lim\limits_{x \rightarrow x_0} \dfrac{f(x)-f\left(x_0\right)}{x-x_0}$}
		{$f'\left(x_0\right)=\lim\limits_{x  \rightarrow x_0} \dfrac{f(x)+f\left(x_0\right)}{x-x_0}$}
		{$f'\left(x_0\right)=\lim\limits_{x \rightarrow x_0} \dfrac{f(x)-f\left(x_0\right)}{x+x_0}$} 
		{$f'\left(x_0\right)=\lim\limits_{x  \rightarrow x_0} \dfrac{f(x)+f\left(x_0\right)}{x+x_0}$}
	 \loigiai{
		Theo định nghĩa đạo hàm  ta có
		$$ 
		f'\left(x_0\right)=\lim\limits_{x \rightarrow x_0} \dfrac{f(x)-f\left(x_0\right)}{x-x_0}. 
		 $$
	 }
\end{ex}

% nguồn: Ok_TracNghiemLop11_Lan1
\begin{ex}%[1D7H1-1]%[ma:0cde7f43]
	Tính đạo hàm của hàm  số $f(x)=\dfrac{x^2+|x+1|}{x}$ tại $x =1$.
	\choice 
		{$ 2  $}
		{$ 0 $} 
		{$ 3 $} 
		{\True Đáp án  khác}
	\loigiai{ 
		Ta có  hàm số liên tục tại $x=1$ và
		 $$
		 f(-1)=-1; \dfrac{f(x)-f(-1)}{x+1}=\dfrac{x^2+x+|x+1|}{x(x+1)}.
		 $$
		 Nên
		\begin{align*}
			 &\lim\limits_{x \rightarrow-1^{+}} \dfrac{f(x)-f(-1)}{x+1}
			=\lim\limits_{x \rightarrow-1^{+}}  \dfrac{x^2+2 x+1}{x(x+1)}
			=\lim\limits_{x  \rightarrow-1^{+}} \dfrac{(x+1)^2}{x(x+1)}
			=\lim\limits_{x \rightarrow-1^{+}}  \dfrac{x+1}{x}
			 =0; \\
			&\lim\limits_{x \rightarrow 1^{-}}  \dfrac{f(x)-f(-1)}{x+1}
			 =\lim\limits_{x \rightarrow 1^{-}} \dfrac{x^2-1}{x(x+1)}
			 =\lim\limits_{x \rightarrow-1^{-}} \dfrac{(x-1) \cdot(x+1)}{x(x+1)}
			=\lim\limits_{x  \rightarrow-1^{-}} \dfrac{x-1}{x}
			=2. 
		\end{align*}
		Do  đó $\lim\limits_{x \rightarrow-1^{+}} \dfrac{f(x)-f(-1)}{x+1} \neq \lim\limits_{x \rightarrow-1^{-}} \dfrac{f(x)-f(-1)}{x+1}$.
		 \\
		Vậy hàm số không có  đạo hàm tại điểm $x =-1$.
		 \\
		Nhận xét: Hàm số $y=f(x)$ có đạo hàm tại $x=x_0$ thì phải liên tục tại  điểm đó.
	 }
\end{ex}

% nguồn: Ok_TracNghiemLop11_Lan1
\begin{ex}%[1D7H1-1]%[ma:e4a4d2e0]
	Cho hàm số $f(x)=\heva{&(x-1)^2, &&x \geq 0 \\ &-x^2, &&x<0}$. $ f(x) $ có đạo hàm tại  điểm $x_0=0$ là
	\choice 
		 {$f'(0)=0$}
		 {$f'(0)=1$}
		{$f'(0)=-2$} 
		{\True Không tồn tại} 
	\loigiai{ 
		 Ta có
		\begin{align*}
			 &f(0)=1;\\
			&\lim\limits_{x \rightarrow  0^{+}} f(x)=\lim\limits_{x \rightarrow 0^{+}}(x-1)^2=1;\\
			&\lim\limits_{x \rightarrow 0^{-}}  f(x)=\lim\limits_{x \rightarrow 0^{-}}\left(-x^2\right)=0.
		\end{align*}
		Do $f(0)=\lim\limits_{x \rightarrow 0^{+}} f(x) \neq \lim\limits_{x \rightarrow 0^{-}} f(x)$ nên hàm số không liên tục tại $x_0=0$. 
		 \\
		Vậy hàm số không  có đạo hàm tại $x_0=0$.
	} 
\end{ex}

% nguồn: Ok_TracNghiemLop11_Lan1
\begin{ex}%[1D7H1-1]%[ma:03399e4d]
	Cho hàm số $y=f(x)=\heva{&x^2+1,  && x \geq 1 \\ &2 x, && x<1}$. Mệnh đề \textbf{sai} là
	\choice 
		 {$f'(1)=2$}
		 {\True $f$ không có đạo hàm tại $x_0=1$}
		{$f'(0)=2$} 
		 {$f'(2)=4$}
	 \loigiai{
		 Ta có
		\begin{align*}
			&\lim\limits_{x \rightarrow 1^{-}} \dfrac{f(x)-f(1)}{x-1}=\lim\limits_{x \rightarrow  1^{-}} \dfrac{2 x-2}{x-1}=\lim\limits_{x \rightarrow 1^{-}} 2=2; \\
			&\lim\limits_{x \rightarrow  1^{+}} \dfrac{f(x)-f(1)}{x-1}=\lim\limits_{x \rightarrow 1^{+}} \dfrac{x^2+1-2}{x-1}=\lim\limits_{x \rightarrow 1^{+}}(x+1)=2 .
		\end{align*}
		Suy ra hàm số có đạo  hàm tại $x_0=1$ là $f'(1)=2$.
	 }
\end{ex}

% nguồn: Ok_TracNghiemLop11_Lan1
\begin{ex}%[1D7H1-1]%[ma:b2aa7900]
	Cho hàm số $f(x)=\heva{&\dfrac{3-\sqrt{4-x}}{4}, &&x \neq 0 \\ &\dfrac{1}{4}, &&x=0}$. Khi đó $f'(0)$  là kết quả nào sau đây?
	 \choice
		{$\dfrac{1}{4}$} 
		{\True $\dfrac{1}{16}$} 
		{$\dfrac{1}{32}$} 
		 {Không tồn tại}
	 \loigiai{
		Với $x \neq 0$ xét 
		\begin{align*}
			\lim\limits_{x \rightarrow  0} \dfrac{f(x)-f(0)}{x-0}
			&=\lim\limits_{x \rightarrow 0} \dfrac{\dfrac{3-\sqrt{4-x}}{4}-\dfrac{1}{4}}{x}  \\
			&=\lim\limits_{x \rightarrow 0} \dfrac{2-\sqrt{4-x}}{4 x}=\lim\limits_{x \rightarrow  0} \dfrac{4-(4-x)}{4 x(2+\sqrt{4-x})} \\
			&=\lim\limits_{x \rightarrow 0} \dfrac{1}{4(2+\sqrt{4-x})}  \\
			&=\dfrac{1}{4(2+\sqrt{4-0})} \\ 
			&=\dfrac{1}{16}. 
		\end{align*}
		Vậy $  f'(0)=\dfrac{1}{16} $
	} 
\end{ex}

% nguồn: Ok_TracNghiemLop11_Lan1
\begin{ex}%[1D7H1-1]%[ma:2482ee4f]
	Cho hàm số $y=\heva{&x^2+a x+b, &&x \geq 2 \\ &x^3-x^2-8 x+10, &&x<2}$. Biết hàm số có đạo hàm tại điểm $x=2$. Giá trị của $a^2+b^2$ bằng 
	 \choice
		{\True $ 20  $}
		{$ 17 $} 
		{$ 18  $}
		 {$ 25 $}
	 \loigiai{
		 Ta có
		 $$
		y=\heva{&x^2+a x+b,  &&x \geq 2 \\ &x^3-x^2-8 x+10, &&x<2.}
		$$ 
		Suy  ra
		$$ 
		y'= \heva{&2 x+a, &&x  \geq 2 \\ &3 x^2-2 x-8 &&x<2.}
		 $$
		Hàm  số có đạo hàm tại điểm $x=2$ nên
		$$ 
		4+a=0  \quad \text{hay} \quad a=-4
		$$ 
		Mặt khác hàm  số có đạo hàm tại điểm $x=2$ thì hàm số liên tục tại điểm $x=2$. Suy ra
		 $$
		\lim\limits_{x \rightarrow  2^{+}} f(x)
		=\lim\limits_{x \rightarrow  2^{-}} f(x)=f(2).
		 $$
		Do  đó
		$$ 
		4+2 a+b=-2 \quad \text{hay} \quad b=2. 
		 $$
		Vậy  $a^2+b^2=20$.
	 }
\end{ex}

% nguồn: Ok_TracNghiemLop11_Lan1
\begin{ex}%[1D7H1-1]%[ma:4fc200d0]
	Cho hàm số $f(x)$ xác định bởi $f(x)=\heva{&\dfrac{\sqrt{4 x^2+1}-1}{x}, &&x \neq  0 \\ &0, &&x=0}$. Giá trị $f'(0)$ bằng
	 \choice
		{\True  $ 2 $}
		{$  0 $}
		{$\dfrac{1}{2}$} 
		{Không tồn tại} 
	 \loigiai{
		Tập xác định  $\mathscr{D} = \mathbb{R}$.
		 \\
		 Ta có
		 $$
		\lim\limits_{x \rightarrow 0}  \dfrac{f(x)-f(0)}{x-0}
		 =\lim\limits_{x \rightarrow 0} \dfrac{\sqrt{4 x^2+1}-1}{x^2}
		=\lim\limits_{x \rightarrow 0} \dfrac{4 x^2}{x^2\left(\sqrt{4  x^2+1}+1\right)}
		=\lim\limits_{x  \rightarrow 0} \dfrac{4}{\sqrt{4 x^2+1}+1}
		 =2.
		 $$
		 Vậy $f'(0)=2$.
	 }
\end{ex}

% nguồn: Ok_TracNghiemLop11_Lan1
\begin{ex}%[1D7H1-1]%[ma:3523f54e]
	Cho hàm số $f(x)=\heva{&\dfrac{x^2}{2}, &&x \leq 1 \\ &a x+b, &&x>1}$. Với giá trị nào sau  đây của $a$, $b$ thì hàm số có đạo hàm tại $x =1$?
	\choice 
		 {\True $a=1; b=-\dfrac{1}{2}$}
		{$a=\dfrac{1}{2}; b=\dfrac{1}{2}$} 
		 {$a=\dfrac{1}{2}; b=-\dfrac{1}{2}$}
		 {$a=1; b=\dfrac{1}{2}$}
	\loigiai{ 
		Hàm số liên tục tại $x=1$ nên ta có 
		\begin{align*}
			\lim\limits_{x \rightarrow 1^{+}}  f(x)&=\lim\limits_{x \rightarrow 1^{-}} f(x) \\
			\dfrac{1}{2}&=a+b. 
		\end{align*}
		Hàm số có đạo hàm tại  $x =1$ nên giới hạn 2 bên của $\dfrac{f(x)-f(1)}{x-1}$ bằng nhau
		\begin{align*}
			\lim\limits_{x \rightarrow  1^{+}} \dfrac{f(x)-f(1)}{x-1}
			&=\lim\limits_{x \rightarrow 1^{+}} \dfrac{a  x+b-(a .1+b)}{x-1} \\
			&=\lim\limits_{x \rightarrow 1^{+}} \dfrac{a(x-1)}{x-1} \\ 
			&=\lim\limits_{x  \rightarrow 1^{+}} a \\
			 &=a
			 \qquad\left(\text {Vì } f(1)=\dfrac{1}{2}=a+b\right)
		\end{align*}
		và 
		$$ 
		\lim\limits_{x \rightarrow  1^{-}} \dfrac{f(x)-f(1)}{x-1}
		=\lim\limits_{x \rightarrow 1^{-}} \dfrac{\dfrac{x^2}{2}-\dfrac{1}{2}}{x-1} 
		 =\lim\limits_{x \rightarrow 1^{-}} \dfrac{(x+1)(x-1)}{2(x-1)}
		 =\lim\limits_{x \rightarrow 1^{-}} \dfrac{(x+1)}{2}=1.
		 $$
		Vậy $a=1$  và $b=-\dfrac{1}{2}$.
	} 
\end{ex}

% nguồn: Ok_TracNghiemLop11_Lan1
\begin{ex}%[1D7H1-1]%[ma:418aa966]
	Cho hàm số $f(x)=\heva{&a x^2+b x+1, &&x \geq 0  \\ &a x-b-1, &&x<0}$. Khi hàm số $f(x)$ có đạo hàm tại $x_0=0$. Hãy tính $T=a+2 b$.
	 \choice
		{$T=-4$} 
		 {$T=0$}
		{\True $T=-6$} 
		 {$T=4$}
	\loigiai{ 
		Ta  có $f(0)=1$.
		\begin{align*}
			&\lim\limits_{x \rightarrow 0^{+}} f(x)=\lim\limits_{x \rightarrow 0^{+}}\left(a  x^2+b x+1\right)=1;\\
			&\lim\limits_{x \rightarrow 0^{-}}  f(x)=\lim\limits_{x \rightarrow 0^{-}}(a x-b-1)=-b-1.
		\end{align*}
		Để hàm  số có đạo hàm tại $x_0=0$ thì hàm số phải liên tục tại $x_0=0$ nên $$
		f(0)=\lim\limits_{x \rightarrow 0^{+}} f(x)=\lim\limits_{x  \rightarrow 0^{-}} f(x).
		$$ 
		Suy ra  $-b-1=1 $ hay $ b=-2$.
		\\ 
		Khi đó $f(x)=\heva{&a x^2-2 x+1, &&x \geq  0 \\ &a x+1, &&x<0}$.
		 \\
		 Xét
		\begin{align*}
			&\lim\limits_{x  \rightarrow 0^{+}} \dfrac{f(x)-f(0)}{x}=\lim\limits_{x \rightarrow 0^{+}} \dfrac{a x^2-2 x+1-1}{x}=\lim\limits_{x \rightarrow 0^{+}}(a x-2)=-2; \\
			&\lim\limits_{x \rightarrow 0^{-}} \dfrac{f(x)-f(0)}{x}=\lim\limits_{x \rightarrow 0^{-}} \dfrac{a x+1-1}{x}=\lim\limits_{x \rightarrow 0^{-}}(a)=a 
		\end{align*}
		 Hàm số có đạo hàm tại $x_0=0$ thì $a=-2$.
		 \\
		Vậy với $a=-2$,  $b=-2$ thì hàm số có đạo hàm tại $x_0=0$ khi đó $T=a+2b=-6$.
	 }
\end{ex}

% nguồn: Ok_TracNghiemLop11_Lan1
\begin{ex}%[1D7H1-1]%[ma:bd1f68c9]
	Cho hàm số $y=f(x)$  có đạo hàm tại điểm $x_0=2$. Tìm $\lim\limits_{x \rightarrow 2} \dfrac{2 f(x)-x f(2)}{x-2}$.
	 \choice
		{$  0 $}
		{$f'(2)$} 
		 {\True $2 f'(2)-f(2)$}
		{$f(2)-2 f'(2)$} 
	\loigiai{ 
		Do hàm số $y=f(x)$ có đạo hàm tại  điểm $x_0=2$ suy ra
		 $$
		\lim\limits_{x \rightarrow  2} \dfrac{f(x)-f(2)}{x-2}=f'(2).
		$$ 
		 Ta có
		\begin{align*}
			 L &=\lim\limits_{x \rightarrow 2} \dfrac{2 f(x)-x f(2)}{x-2} \\
			 &=\lim\limits_{x \rightarrow 2} \dfrac{2 f(x)-2 f(2)+2 f(2)-x f(2)}{x-2} \\
			&=\lim\limits_{x \rightarrow 2} \dfrac{2(f(x)-f(2))}{x-2}-\lim\limits_{x \rightarrow  2} \dfrac{f(2)(x-2)}{x-2} \\
			 &=2 f'(2)-f(2).
		\end{align*}
	 }
\end{ex}

% nguồn: Ok_TracNghiemLop11_Lan1
\begin{ex}%[1D7H1-1]%[ma:0b253228]
	Cho hàm số $f(x)=\heva{&\dfrac{\sqrt{3 x+1}-2 x}{x-1}, &&x \neq 1 \\ &\dfrac{-5}{4}  &&x=1}$. Tính $f'(1)$.
	\choice 
		 {Không tồn tại}
		{$ 0  $}
		 {$-\dfrac{7}{50}$}
		{\True $-\dfrac{9}{64}$} 
	 \loigiai{
		Ta có 
		$$ 
		\lim\limits_{x \rightarrow  1} f(x)
		 =\lim\limits_{x \rightarrow 1} \dfrac{\sqrt{3 x+1}-2 x}{x-1}
		=\lim\limits_{x \rightarrow 1} \dfrac{3 x+1-4 x^2}{(x-1)(\sqrt{3 x+1}+2 x)} 
		=\lim\limits_{x  \rightarrow 1} \dfrac{-4 x-1}{(\sqrt{3 x+1}+2 x)}
		 =\dfrac{-5}{4}=f(1)
		 $$
		Suy ra hàm số liên tục lại  $x=1$.
		\\ 
		 Mặt khác
		\begin{align*}
			f'(1) & =\lim\limits_{x \rightarrow  1} \dfrac{f(x)-f(1)}{x-1} \\
			&=\lim\limits_{x \rightarrow 1} \dfrac{\dfrac{\sqrt{3  x+1}-2 x}{x-1}+\dfrac{5}{4}}{x-1}\\
			&=\lim\limits_{x \rightarrow 1} \dfrac{4 \sqrt{3  x+1}-3 x-5}{4(x-1)^2} \\
			& =\lim\limits_{x \rightarrow  1} \dfrac{16(3 x+1)-(3 x+5)^2}{4(x-1)^2(4 \sqrt{3 x+1}+3 x+5)}\\
			&=\lim\limits_{x \rightarrow 1} \dfrac{-9}{4(4 \sqrt{3 x+1}+3  x+5)}\\
			 &=-\dfrac{9}{64}.
		\end{align*}
	 }
\end{ex}

% nguồn: Ok_TracNghiemLop11_Lan1
\begin{ex}%[1D7H1-1]%[ma:5bd91036]
	Cho hàm số $f(x)=\heva{&a x^2+b x, &&x \geq 1 \\ &2 x-1, &&x<1}$. Để hàm số đã  cho có đạo hàm tại $x=1$ thì $2 a+b$ bằng
	\choice 
		 {\True $ 2 $}
		 {$ 5 $}
		{$ -2  $}
		{$  -5 $}
	\loigiai{ 
		Ta  có
		\begin{align*}
			 \lim\limits_{x \rightarrow 1^{-}} \dfrac{f(x)-f(1)}{x-1}
			=\lim\limits_{x \rightarrow 1^{-}} \dfrac{2  x-1-1}{x-1}=2;
		\end{align*}
		 và
		\begin{align*}
			\lim\limits_{x \rightarrow 1^{+}} \dfrac{f(x)-f(1)}{x-1} 
			&=\lim\limits_{x \rightarrow 1^{+}} \dfrac{a x^2+b x-a-b}{x-1}\\ 
			&=\lim\limits_{x \rightarrow 1^{+}}  \dfrac{a\left(x^2-1\right)+b(x-1)}{x-1}\\
			&=\lim\limits_{x \rightarrow 1^{+}}  \dfrac{(x-1)[a(x+1)+b]}{x-1} \\
			&=\lim\limits_{x \rightarrow  1^{+}}[a(x+1)+b]\\
			&=2 a+b. 
		\end{align*}
		Theo yêu cầu bài  toán
		\begin{align*}
			\lim\limits_{x \rightarrow 1^{-}} \dfrac{f(x)-f(1)}{x-1}&=\lim\limits_{x \rightarrow 1^{+}}  \dfrac{f(x)-f(1)}{x-1} \\
			 2 a+b&=2.
		\end{align*}
	 }
\end{ex}

% nguồn: Ok_TracNghiemLop11_Lan1
\begin{ex}%[1D7H1-1]%[ma:6b1f7460]
	Tính đạo hàm của hàm số $y=2  x^2+x+1$ tại điểm $x=2$
	 \choice
		 {\True $ 9 $}
		{$ 4  $}
		{$ 7 $} 
		{$ 6 $} 
	\loigiai{ 
		Cho $x_0=2$ một số gia $\Delta x$. 
		Khi đó hàm số nhận một số gia tương ứng 
		\begin{align*}
			 \Delta y
			&  =f\left(x_0+\Delta x\right)-f\left(x_0\right) \\
			&  =2(2+\Delta x)^2+(2+\Delta x)+1-\left(2.2^2+2+1\right) \\
			& =8+8 \Delta x+2(\Delta x)^2+2+\Delta x+1-11 \\ 
			& =\Delta x(9+2 \Delta  x).
		\end{align*}
		Ta có 
		 $$
		f'(2)=\lim\limits_{\Delta x \rightarrow 0} \dfrac{\Delta y}{\Delta  x}
		=\lim\limits_{\Delta x \rightarrow 0} \dfrac{\Delta x(9+2 \Delta  x)}{\Delta x}
		=\lim\limits_{\Delta x \rightarrow 0}(9+2  \Delta x)
		=9. 
		 $$
	} 
\end{ex}

% nguồn: Ok_Toan11_TracNghiem_Lan2
\begin{ex}%[1D7V1-1]%[ma:a19e5729]
	Cho hàm số $ f(x)=\heva{&a x^{2}+b x+1\text{ khi } x \geq 0 \\& a x-b-1 \text{ khi } x<0}$. Khi hàm số $ f(x) $ có  đạo hàm tại $ x_{0}=0 $ thì $a+2 b $ bằng
	 \choice
		{$T=-4$} 
		{$T=0$} 
		{\True $T=-6$} 
		 {$T=4$}
	 \loigiai{
		Ta có $ f(0)=1  $ và
		\begin{itemize}
			\item $\lim\limits_{x \rightarrow  0^{+}} f(x)=\lim\limits_{x \rightarrow 0^{+}}\left(a x^{2}+b x+1\right)=1$.
			\item $\lim\limits_{x \rightarrow  0^{-}} f(x)=\lim\limits_{x \rightarrow 0^{-}}(a x-b-1)=-b-1$.
		\end{itemize}
		Để hàm số có đạo hàm tại $ x_{0}=0 $ thì hàm số phải liên tục  tại $ x_{0}=0 $ nên
		$$ f(0)=\lim\limits_{x  \rightarrow 0^{+}} f(x)=\lim\limits_{x \rightarrow 0^{-}} f(x)\Leftrightarrow -b-1=1\Leftrightarrow b=-2. $$.
		Khi đó $ f(x)=\heva{&a x^{2}-2 x+1\text{ khi }  x \geq 0 \\& a x+1\text{ khi } x<0.}$\\
		 Xét
		\begin{itemize}
			\item $ \lim\limits_{x \rightarrow 0^{+}}  \dfrac{f(x)-f(0)}{x}=\lim\limits_{x \rightarrow 0^{+}} \dfrac{a x^{2}-2 x+1-1}{x}=\lim\limits_{x \rightarrow 0^{+}}(a x-2)=-2 $.
			 \item $ \lim\limits_{x \rightarrow 0^{-}} \dfrac{f(x)-f(0)}{x}=\lim\limits_{x \rightarrow 0^{-}} \dfrac{a x+1-1}{x}=\lim\limits_{x \rightarrow 0^{-}}(a)=a $.
		\end{itemize}
		Hàm số có  đạo hàm tại $ x_{0}=0 $ thì $ a=-2 $.\\
		Vậy với $ a=-2$, $b=-2 $ thì  hàm số có đạo hàm tại $ x_{0}=0 $ khi đó $ T=a+2b=-6 $.
	 }
\end{ex}

% nguồn: Ok_Toan11_TracNghiem_Lan2
\begin{ex}%[1D7V1-1]%[ma:4c99875b]
	Cho hàm số $f(x)=\heva{& \dfrac{3-\sqrt{4-x}}{4}&& \text{khi  } x\ne 0\\&\dfrac{1}{4}&& \text{khi } x=0}$. Tính $f'(0)$.
	\choice 
		{Không  tồn tại}
		 {\True $f'(0)=\dfrac{1}{16}$}
		 {$f'(0)=\dfrac{1}{4}$}
		 {$f'(0)=\dfrac{1}{32}$}
	\loigiai{Ta có \begin{align*}
		\lim\limits_{x\to 0}\dfrac{f(x)-f(0)}{x-0}&=\lim\limits_{x\to 0}\dfrac{\dfrac{3-\sqrt{4-x}}{4}-\dfrac{1}{4}}{x}=\lim\limits_{x\to 0}\dfrac{2-\sqrt{4-x}}{4x}=\lim\limits_{x\to 0}\dfrac{4-(4-x)}{4x\left(2+\sqrt{4-x}\right)}\\ 
		&=\lim\limits_{x\to 0}\dfrac{1}{4\left(2+\sqrt{4-x}\right)}=\dfrac{1}{16}. 
	\end{align*}
	 Vậy $f'(0)=\dfrac{1}{16}$.
} 
\end{ex}

% nguồn: Ok_TracNghiemLop11_Lan1
\begin{ex}%[1D7Y1-1]%[ma:4a804c0c]
	Chọn  khẳng định đúng trong các khẳng định sau
	\choice 
		{Hàm số $y=f(x)$ liên tục tại điểm $x_0$ thì có  đạo hàm tại điểm đó}
		{\True Hàm số $y=f(x)$ có đạo hàm tại điểm  $x_0$ thì liên tục tại điểm đó}
		{Hàm số $y=f(x)$ xác định tại điểm $x_0$ thì có đạo hàm  tại điểm đó}
		{Hàm số $y=f(x)$ luôn có đạo hàm tại mọi điểm thuộc tập  xác định của nó}
	\loigiai{Nếu hàm số $y=f(x)$ có đạo hàm tại điểm $x_0$ thì liên tục  tại điểm đó.
	} 
\end{ex}

% nguồn: Ok_TracNghiemLop11_Lan1
\begin{ex}%[1D7K1-1]%[ma:86588c1f]
	Cho hàm số $ y=f(x)= \begin{cases} \sqrt{4x+1} & \text{khi } x>0 \\ x+1 & \text{khi } x \leq 0
\end{cases}$. Khẳng định nào là đúng về đạo hàm của hàm số $f(x)$ tại $x=0$?
 \choice
	{$f'(0)  = - 1$}
	 {$f'(0) = 0$}
	{$f'(0) = 1$} 
	{\True  Không tồn tại}
 \loigiai{
	$\displaystyle \lim_{x \rightarrow 0^+} \dfrac{f(x) - f(0)}{x-0} 
	 = \lim_{x \rightarrow 0^+} \dfrac{\sqrt{4x+1} - 1}{x}
	=\lim_{x \rightarrow 0^+} \dfrac{4x  + 1 - 1}{x \left(\sqrt{4x+1} + 1\right)} = 2$;\\
	$\displaystyle \lim_{x \rightarrow 0^-} \dfrac{f(x) -  f(0)}{x-0}
	= \lim_{x  \rightarrow 0^-} \dfrac{x+1 - 1}{x-0} = 1$.\\
	Do $\displaystyle \lim_{x \rightarrow 0^+} \dfrac{f(x) - f(0)}{x-0} \ne \lim_{x \rightarrow 0^-} \dfrac{f(x) - f(0)}{x-0}$ nên $\displaystyle \lim_{x  \rightarrow 0} \dfrac{f(x) - f(0)}{x-0}$ không tồn tại.
} 
\end{ex}


%%%%%%%%%%%%%%%% ĐÚNG/SAI %%%%%%%%%%%%%%%%

% nguồn: Ok_dungsai_lop11_Lan1
\begin{ex}%[1D7N1-1]%[ma:b77a893c]
 Dùng định nghĩa để tính đạo hàm  của hàm số $y=f(x)=x^2+2x$ tại điểm $x_0=1$. Khi đó
 \choiceTF 
  {\True $f'(1)=\displaystyle\lim_{x \to  1}\dfrac{f(x)-f(1)}{x-1}$}
  {\True $f'(1)=\displaystyle\lim_{x  \to 1}\dfrac{x^2+2x-3}{x-1}$}
   {$f'(1)=\displaystyle\lim_{x\to 1}\left( x+4\right)$ }
  {$f'(1)=a \Rightarrow  a>5$}
 \loigiai{Theo định nghĩa ta có  \\
  $f'(1)=\displaystyle\lim_{x \to 1}\dfrac{f(x)-f(1)}{x-1}=\displaystyle\lim_{x \to  1}\dfrac{x^2+2x-3}{x-1}=\displaystyle\lim_{x \to 1}\dfrac{(x-1)(x+3)}{x-1}=\displaystyle\lim_{x \to 1}(x+3)=4$.
 } 
\end{ex}

% nguồn: Ok_dungsai_lop11_Lan1
\begin{ex}%[1D7N1-1]%[ma:0a3fcab1]
 Cho hàm số $y=f(x)=x^2+1$ và  hàm số $y=g(x)=\dfrac{x+1}{x-1}$. Xét tính đúng sai của các mệnh đề sau
 \choiceTF 
  {Hàm số $y=f(x)$ và  $y=g(x)$ có đạo hàm trên $\mathbb{R}$}
   {\True Hàm số $y=f(x)$ có đạo hàm trên $\mathbb{R}$, hàm số$y=g(x)$ có đạo hàm trên $\mathbb{R}\backslash\left\{1\right\}$}
  {\True $f'(x)=2x, \forall  x\in\mathbb{R}$.}
  {\True $g'(2)=-2$} 
 \loigiai{ 
  \begin{itemchoice}
   \itemch Sai vì hàm số $y=g(x)=\dfrac{x+1}{x-1} $có tập xác định là  $\mathbb{R}\backslash\left\{1\right\}$ nên có đạo hàm trên $\mathbb{R}\backslash\left\{1\right\}$.
   \itemch Đúng vì Hàm số $y=f(x)=x^2+1$ có tập xác định  là $\mathbb{R}$ nên có đạo hàm trên $\mathbb{R}$.\\
   Hàm số $y=g(x)=\dfrac{x+1}{x-1} $có tập xác định là $\mathbb{R}\backslash\left\{1\right\}$ nên có đạo hàm trên  $\mathbb{R}\backslash\left\{1\right\}$.
   \itemch  Đúng vì $f'(x)=\left(x^2+1\right)'=2x$.
   \itemch Đúng  vì $g(x)=\dfrac{-2}{(x-1)^2}\Rightarrow g'(2)=-2$.
  \end{itemchoice}
  }
\end{ex}

% nguồn: Ok_dungsai_lop11_Lan1
\begin{ex}%[1D7N1-1]%[ma:27a86d1a]
 Cho hàm số $y=f(x)$ có đạo hàm tại $x_0$ là $f'(x_0)$.  Hãy chỉ ra tính đúng, sai của các khẳng định sau?
 \choiceTF 
  {\True  $f'(x_0)=\lim \limits_{x\to x_0} \dfrac{f(x)-f(x_0)}{x-x_0}$}
  {\True  $f'(x_0)=\lim \limits_{\Delta x\to 0} \dfrac{f(x_0+\Delta x)-f(x_0)}{\Delta x}$}
   {\True $f'(x_0)=\lim \limits_{h\to 0} \dfrac{f(x_0+h)-f(x_0)}{h}$}
  {$f'(x_0)=\lim \limits_{x\to x_0} \dfrac{f(x+x_0)-f(x_0)}{x-x_0}$} 
  \loigiai{
  Hàm số $y=f(x)$ có đạo hàm tại $x_0$ là $f'(x_0) \Rightarrow f'(x_0)=\lim \limits_{x\to x_0}  \dfrac{f(x)-f(x_0)}{x-x_0}$.\\
   Đặt $h=\Delta x=x-x_0\Rightarrow f'(x_0)=\lim \limits_{\Delta x\to 0} \dfrac{f(x_0+\Delta x)-f(x_0)}{\Delta x}=\lim \limits_{h\to 0} \dfrac{f(x_0+h)-f(x_0)}{h}$.\\
 Vậy  chỉ có khẳng định $f'(x_0)=\lim \limits_{x\to x_0} \dfrac{f(x+x_0)-f(x_0)}{x-x_0}$ sai, các khẳng định còn lại đúng.}
\end{ex}

% nguồn: Ok_dungsai_lop11_Lan1
\begin{ex}%[1D7H1-1]%[ma:30edde38]
 Dùng định nghĩa để tính đạo hàm của hàm số $y=f(x)=x^2+2 x$ tại điểm $x_0=1$. Khi đó, các mệnh đề sau đúng  hay sai?
  \choiceTF
  {\True  $f'(1)=\lim\limits_{x \rightarrow 1} \dfrac{f(x)-f(1)}{x-1}$}
  {\True $f'(1)=\lim\limits_{x  \rightarrow 1} \dfrac{x^2+2 x-3}{x-1}$}
  {$f'(1)=\lim\limits_{x \rightarrow 1}(x+4)$} 
  {$f'(1)=a \Rightarrow a>5$} 
 \loigiai{ 
  Ta  có
   $$
   f'(1)=\lim\limits_{x \rightarrow 1} \dfrac{f(x)-f(1)}{x-1}
   =\lim\limits_{x \rightarrow 1} \dfrac{x^2+2 x-3}{x-1}
  =\lim\limits_{x \rightarrow 1}  \dfrac{(x-1)(x+3)}{x-1}
  =\lim\limits_{x  \rightarrow 1}(x+3)
  =4. 
  $$ 
   Vậy $f'(1)=4<5$.
   \\
  Dựa vào biến đổi  trên, ta thấy
  \begin{itemchoice}
   \itemch $ f'(1)=\lim\limits_{x  \rightarrow 1} \dfrac{f(x)-f(1)}{x-1} $.
   \itemch $ f'(1)=\lim\limits_{x \rightarrow  1} \dfrac{x^2+2 x-3}{x-1} $.
   \itemch $f'(1)=\lim\limits_{x \rightarrow  1}(x+3)$.
   \itemch $f'(1)=a \Rightarrow  a=4$.
  \end{itemchoice}
 } 
\end{ex}

% nguồn: Ok_dungsai_lop11_Lan1
\begin{ex}%[1D7H1-1]%[ma:87eb66c4]
 Cho hàm số $f(x)=\heva{& x^2-1 \text{ khi  } x\geqslant 0 \\& -x^2 \text{ khi } x<0 } $. Xét tính đúng, sai của các khẳng định sau
  \choiceTF
   {\True Hàm số không liên tục tại $x=0$}
  {Hàm số không có đạo hàm tại  $x=-1$}
  {\True Hàm số  có đạo hàm tại $x=2$}
  {Hàm số không  liên tục tại $x=2$}
 \loigiai{ 
  \begin{itemize}
   \item Xét các giới hạn $\heva{&\lim  \limits_{x\to 0^{+}} f(x)=\lim \limits_{x\to 0^{+}} (x^2-1)=-1\\&\lim \limits_{x\to 0^{-}} f(x)=\lim \limits_{x\to 0^{-}} (-x^2)=0.}$\\
   Vì $\lim \limits_{x\to 0^{+}}  f(x)\ne \lim \limits_{x\to 0^{-}} f(x)$ nên hàm số không liên tục tại $x=0$.\\
   \item Với $x<0$, $f(x)=-x^2\Rightarrow  $ do đó với $x=-1$, hàm số có đạo hàm tại điểm đó.
   \item Với $x\geq 0$, $f(x)=x^2-1\Rightarrow $. Do đó tại điểm $x=2$, hàm số liên tục và  có đạo hàm tại điểm đó.
 \end{itemize}}
\end{ex}

% nguồn: Ok_dungsai_lop11_Lan1
\begin{ex}%[1D7H1-1]%[ma:d2c3adb0]
 Dùng định nghĩa để tính đạo hàm của hàm số $f(x)=2 x^3$. Khi đó, các mệnh đề sau đúng  hay sai?
 \choiceTF 
  {\True Với bất kì $x_0 \colon f'\left(x_0\right)=\lim\limits_{x \rightarrow x_0}  \dfrac{f(x)-f\left(x_0\right)}{x-x_0}$}
  {$f'(1)=-6$} 
   {\True $f'(0)=0$}
  {\True $f'(2)=24$} 
  \loigiai{
  Với bất kì $x_0$,  ta có
  \begin{align*}
   f'\left(x_0\right)  &=\lim\limits_{x \rightarrow x_0} \dfrac{f(x)-f\left(x_0\right)}{x-x_0} \\
    &=\lim\limits_{x \rightarrow x_0} \dfrac{2 x^3-2 x_0^3}{x-x_0} \\
   &=\lim\limits_{x \rightarrow x_0} \dfrac{2\left(x-x_0\right)\left(x^2+x  \cdot x_0+x_0^2\right)}{x-x_0} \\
   &=\lim\limits_{x \rightarrow  x_0} \dfrac{2\left(x^2+x \cdot x_0+x_0^2\right)}{1} \\
   &=6 x_0^2 . 
  \end{align*}
  Khi  đó $f'(x)=\left(2 x^3\right)'=6 x^2$ trên $\mathbb{R}$.
   \\
  Ta thấy $  f'(1)=6 $, $ f'(0)=0 $ và $ f'(2)=24 $.
  }
\end{ex}

% nguồn: Ok_dungsai_lop11_Lan1
\begin{ex}%[1D7H1-1]%[ma:2f5540df]
 Dùng định nghĩa để tính đạo hàm của hàm số $f(x)=\dfrac{2}{1-x}$ với $x \neq 1$. Khi đó, các  mệnh đề sau đúng hay sai?
  \choiceTF
  {\True Với bất kì $x_0 \neq 1$, ta có  $f'\left(x_0\right)=\lim\limits_{x \rightarrow x_0} \dfrac{2}{(1-x)\left(1-x_0\right)}$}
  {\True  $f'(2)=2$}
   {$f'(3)=\dfrac{1}{3}$}
  {\True  $f'(2)+f'(3)=\dfrac{5}{2}$}
  \loigiai{
  Với bất kì $x_0 \neq 1$,  ta có:
  \begin{align*}
    f'\left(x_0\right)
   =\lim\limits_{x  \rightarrow x_0} \dfrac{f(x)-f\left(x_0\right)}{x-x_0}
    &=\lim\limits_{x \rightarrow x_0} \dfrac{\dfrac{2}{1-x}-\dfrac{2}{1-x_0}}{x-x_0} \\
   &=\lim\limits_{x  \rightarrow x_0} \dfrac{2\left(x-x_0\right)}{(1-x)\left(1-x_0\right)\left(x-x_0\right)} \\
   &=\lim\limits_{x \rightarrow x_0}  \dfrac{2}{(1-x)\left(1-x_0\right)} \\
   &=\dfrac{2}{\left(1-x_0\right)^2}. 
  \end{align*}
  Khi đó  $f'(x)=\left(\dfrac{2}{1-x}\right)'=\dfrac{2}{(1-x)^2}$ trên các khoảng $(-\infty; 1)$ và $(1;+\infty)$.
  Suy  ra
   $$
  f'(2)=2;\, f'(3)=\dfrac{1}{2} 
  \Rightarrow 
   f'(2)+f'(3) = \dfrac{5}{2}.
   $$
  }
\end{ex}

% nguồn: Ok_dungsai_lop11_Lan1
\begin{ex}%[1D7H1-1]%[ma:f2ad6d80]
 Các mệnh đề sau đúng  hay sai?
  \choiceTF
   {\True $y=x^2-x$ tại $x_0=1$ có $f'(1)=1$}
  {$y=\sqrt{x}$  tại $x_0=1$ có $f'(1)=1$}
  {\True $y=\dfrac{1}{x^2+1}$  tại $x_0=0$ có $f'(0)=0$}
  {\True $y=\dfrac{1}{x+1}$ tại $x_0=2$ có  $f'(2)=-\dfrac{1}{9}$}
 \loigiai{ 
  \begin{itemchoice}
   \itemch $f'(1)=\lim\limits_{x  \rightarrow 1} \dfrac{f(x)-f\left(x_0\right)}{x-x_0}=\lim\limits_{x \rightarrow 1} \dfrac{x^2-x-f(1)}{x-1}=\lim\limits_{x \rightarrow 1} \dfrac{x^2-x}{x-1}=\lim\limits_{x \rightarrow 1} x=1$.
   \itemch $f'(1)=\lim\limits_{x  \rightarrow 1} \dfrac{f(x)-f\left(x_0\right)}{x-x_0}=\lim\limits_{x \rightarrow 1} \dfrac{\sqrt{x}-f(1)}{x-1}=\lim\limits_{x \rightarrow 1} \dfrac{\sqrt{x}-1}{x-1}=\lim\limits_{x \rightarrow 1} \dfrac{1}{\sqrt{x}+1}=\dfrac{1}{2}$.
   \itemch $f'(0)=\lim\limits_{x \rightarrow 0} \dfrac{f(x)-f\left(x_0\right)}{x-x_0}=\lim\limits_{x \rightarrow 0} \dfrac{\dfrac{1}{x^2+1}-f(0)}{x}=\lim\limits_{x \rightarrow 0} \dfrac{\dfrac{1}{x^2+1}-1}{x}=\lim\limits_{x \rightarrow 0} \dfrac{-x}{x^2+1}=0$. 
    \itemch $f'(2)=\lim\limits_{x \rightarrow 2} \dfrac{f(x)-f\left(x_0\right)}{x-x_0}=\lim\limits_{x \rightarrow 2} \dfrac{\dfrac{1}{x+1}-f(2)}{x-2}=\lim\limits_{x \rightarrow 2} \dfrac{\dfrac{1}{x+1}-\dfrac{1}{3}}{x-2}=\lim\limits_{x \rightarrow 2} \dfrac{-1}{3(x+1)}=-\dfrac{1}{9}$.
  \end{itemchoice}
  }
\end{ex}

% nguồn: Ok_dungsai_lop11_Lan1
\begin{ex}%[1D7H1-1]%[ma:d36e1749]
 Tính được đạo hàm của mỗi hàm số sau tại các  điểm đã chỉ ra. Khi đó:
 \choiceTF 
  {\True $y=x^2-x$ tại $x_0=1$ có ${{f}^{\prime  }}(1)=1$}
  {$y=\sqrt{x}$ tại $x_0=1$ có ${{f}^{\prime  }}(1)=1$}
  {\True  $y=\dfrac{1}{x^2+1}$ tại $x_0=0$ có ${{f}^{\prime }}(0)=0$}
  {\True $y=\dfrac{1}{x+1}$ tại $x_0=2$ có ${{f}^{\prime }}(2)=-\dfrac{1}{9}$} 
 \loigiai{ 
  \begin{itemize}
   \item $f^{\prime}(1)=\displaystyle\lim _{x  \rightarrow 1} \dfrac{f(x)-f\left(x_0\right)}{x-x_0}=\displaystyle\lim _{x \rightarrow 1} \dfrac{x^2-x-f(1)}{x-1}=\displaystyle\lim _{x \rightarrow 1} \dfrac{x^2-x}{x-1}=\displaystyle\lim _{x \rightarrow 1} x=1$.
   \item $f^{\prime}(1)=\displaystyle\lim _{x \rightarrow  1} \dfrac{f(x)-f\left(x_0\right)}{x-x_0}=\displaystyle\lim _{x \rightarrow 1} \dfrac{\sqrt{x}-f(1)}{x-1}=\displaystyle\lim _{x \rightarrow 1} \dfrac{\sqrt{x}-1}{x-1}=\displaystyle\lim _{x \rightarrow 1} \dfrac{1}{\sqrt{x}+1}=\dfrac{1}{2}$.
   \item $f^{\prime}(0)=\displaystyle\lim _{x \rightarrow 0} \dfrac{f(x)-f\left(x_0\right)}{x-x_0}=\displaystyle\lim  _{x \rightarrow 0} \dfrac{\dfrac{1}{x^2+1}-f(0)}{x}=\displaystyle\lim _{x \rightarrow 0} \dfrac{\dfrac{1}{x^2+1}-1}{x}=\displaystyle\lim _{x \rightarrow 0} \dfrac{-x}{x^2+1}=0$.
   \item $f^{\prime}(2)=\displaystyle\lim _{x \rightarrow 2} \dfrac{f(x)-f\left(x_0\right)}{x-x_0}=\displaystyle\lim _{x \rightarrow 2} \dfrac{\dfrac{1}{x+1}-f(2)}{x-2}=\displaystyle\lim _{x \rightarrow 2} \dfrac{\dfrac{1}{x+1}-\dfrac{1}{3}}{x-2}=\displaystyle\lim  _{x \rightarrow 2} \dfrac{-1}{3(x+1)}=-\dfrac{1}{9}$.
  \end{itemize}
 } 
\end{ex}

% nguồn: Ok_dungsai_lop11_Lan1
\begin{ex}%[1D7H1-1]%[ma:b09caade]
 Cho hàm số $f(x) = \heva{&(m-1)x^2 + m &\text{khi} \ x \geq  2 \\&3x -1 &\text{khi}\ x < 2}(m \text{ là tham số)}$. Khi đó
 \choiceTF 
  {\True Hàm số xác định  với mọi $x\in \mathbb{R}$}
   {$f(2)=5$}
   {Với $m = 2$ thì hàm số liên tục trên $\mathbb{R}$}
  {\True Hàm  số không có đạo hàm tại $x=2$}
 \loigiai{ 
  \begin{itemchoice}
   \itemch Đúng\\ 
   Với $x\geq 2$ thì $f(x)=(m-1)x^2 + m $ là hàm số đa thức nên xác định. Với $x<2$ thì $f(x)=3x-1$ là hàm số đa thức nên xác định. Vậy hàm số xác định với mọi $x\in  \mathbb{R}$.
    \itemch Sai\\
    Ta có $f(2)=4(m-1)+m=5m-4$.
   \itemch  Sai\\
   Để hàm số liên tục trên $\mathbb{R}$ thì chỉ cần hàm số liên tục thêm tại  $x=2$. Do đó
   \[\lim \limits_{x\to 2^+}f(x)=\lim \limits_{x\to 2^-}f(x)=f(2)\Leftrightarrow  4(m-1)+m=5\Leftrightarrow m=\dfrac{9}{5}.\]
    \itemch Đúng\\
   Để hàm số có  đạo hàm tại $x=2$ thì hàm số phải liên tục tại $x=2$. Do đó $m=\dfrac{9}{5}$.\\
   Với $m=\dfrac{9}{5}$ thì $f(x)=\heva{&\tfrac{4}{5}x^2+\tfrac{9}{5}&\text{khi} \ x  \geq 2\\&3x -1 &\text{khi}\ x < 2}$. \\
   Ta có $f'(2^+)=\dfrac{16}{5}$ và $f'(2^-)=3$. Vì $f'(2^+)\ne f'(2^-)$ nên hàm số  không có đạo hàm tại $x=2$.
  \end{itemchoice}
 } 
\end{ex}

% nguồn: Ok_dungsai_lop11_Lan1
\begin{ex}%[1D7H1-1]%[ma:e0aaf604]
 Dùng định nghĩa để tính đạo hàm của  hàm số $f(x)=\dfrac{2}{1-x}$ với $x\ne 1$. Các khẳng định sau là đúng hay sai?
 \choiceTF 
  {\True Với bất kì $x_0\ne  1$, ta có $f'(x)=\lim\limits_{x\to x_0}\dfrac{2}{(1-x)(1-x_0)}$}
  {\True  $f'(2)=2$}
   {$f'(3)=\dfrac{1}{3}$}
  {$f'(2)+f'(3)=\dfrac{3}{2}$} 
 \loigiai{Với $x_0\ne 1$, ta  có
 \begin{itemize}
  \item \begin{align*}
   f'(x_0)=\lim\limits_{x\to x_0}\dfrac{f(x)-f(x_0)}{x-x_0}&=\lim\limits_{x\to x_0}\dfrac{\tfrac{2}{1-x}-\tfrac{2}{1-x_0}}{x-x_0}\\&= \lim\limits_{x\to x_0}\dfrac{2(x-x_0)}{(1-x)(1-x_0)(x-x_0)}\\&=\lim\limits_{x\to x_0}\dfrac{2}{(1-x)(1-x_0)}=\dfrac{2}{(1-x_0)^2}
 \end{align*}
 Vậy $f'(x)=\left(\dfrac{2}{1-x}\right)'=\dfrac{2}{(1-x)^2}$ trên các khoảng $(-\infty;1)$ và  $(1;+\infty)$.
  \item Ta có $f'(x)=\left(\dfrac{2}{1-x}\right)'=\dfrac{2}{(1-x)^2}$ nên $f'(2)=2$.
 \item Ta có $f'(x)=\left(\dfrac{2}{1-x}\right)'=\dfrac{2}{(1-x)^2}$ nên $f'(3)=\dfrac{1}{2}$. 
  \item Ta có $f'(2)+f'(3)=2+\dfrac{1}{2}=\dfrac{5}{2}$.
\end{itemize}
} 
\end{ex}

% nguồn: Ok_dungsai_lop11_Lan1
\begin{ex}%[1D7H1-1]%[ma:a37bba57]
 Cho hàm số $f(x)=\heva{&2x-3&\text{khi }&x<1\\&-\frac{1}{x}&\text{khi }&x\geq 1}$. Xét tính đúng sai của các  mệnh đề sau
  \choiceTF
  {\True Hàm số $f(x)$ liên tục tại  $x=1$}
  {Hàm số  $f(x)$ có đạo hàm tại $x=1$}
  {Hàm số $f(x)$ liên tục tại $x=1$ và hàm số $f(x)$  cũng có đạo hàm tại $x=1$}
  {\True Hàm  số $f(x)$ không có đạo hàm tại $x=1$}
  \loigiai{
  \begin{itemchoice}
   \itemch Đúng vì $\displaystyle\lim\limits_{x\to 1^-}f(x)=\displaystyle\lim\limits_{x\to 1^-}  (2x-3)=-1$ và $\displaystyle\lim\limits_{x\to 1^+} f(x)=\displaystyle\lim\limits_{x\to 1^+} -\dfrac{1}{x}=-1$.\\
   Do đó, hàm số  $f(x)$ liên tục tại $x=1$.
   \itemch Sai vì $\displaystyle\lim\limits_{x\to  1^-} \dfrac{f(x)-f(1)}{x-1}=\displaystyle\lim\limits_{x\to1^-}\dfrac{2(x-1)}{x-1}=\displaystyle\lim\limits_{x\to 1^-}2=2$\\
   và  $\displaystyle\lim\limits_{x\to 1^+}\dfrac{f(x)-f(1)}{x-1}=\displaystyle\lim\limits_{x\to 1^+}\dfrac{-1+x}{x(x-1)}=\displaystyle\lim\limits_{x\to 1^+}\dfrac{1}{x}=1$.\\
   Do  đó, hàm số $f(x)$ không có đạo hàm tại $x=1$.
    \itemch Sai.
   \itemch Đúng. 
  \end{itemchoice}
 } 
\end{ex}

% nguồn: Ok_dungsai_lop11_Lan1
\begin{ex}%[1D7H1-1]%[ma:ca4ae092]
 Dùng định nghĩa để tính đạo hàm của hàm  số $f(x)=\dfrac{x-2}{x+1}$ tại điểm $x_0=0$ ta được $f^{\prime}(0)=a$.
 Khi đó 
 \choiceTF 
  {\True $f^{\prime}(0)=\lim\limits_{x  \to 0} \dfrac{f(x)-f(0)}{x-0}$}
  {$f^{\prime}(0)=\lim\limits_{x  \to 0} \dfrac{4}{x+1}$}
  {\True Phương trình  $3^x=3$ có nghiệm bằng $x=a-2$}
  {$\log _a 9=3$} 
  \loigiai{
   Ta có
  \allowdisplaybreaks 
  \begin{eqnarray*}
   f^{\prime}(0) & =&\lim\limits_{x \rightarrow 0} \dfrac{f(x)-f(0)}{x-0}=\lim\limits_{x \rightarrow 0}  \dfrac{\frac{x-2}{x+1}-(-2)}{x} \\
   & =&\lim\limits_{x \rightarrow  0} \dfrac{3 x}{x(x+1)}=\lim\limits_{x \rightarrow 0} \dfrac{3}{x+1}=3.
  \end{eqnarray*}
  Vậy $f^{\prime}(0)=3$. Do đó $a=3$. Khi đó 
  \begin{itemize}
   \item $f^{\prime}(0) =\lim\limits_{x \rightarrow 0}  \dfrac{f(x)-f(0)}{x-0}$.
   \item $\lim\limits_{x \rightarrow 0} \dfrac{4}{x+1}=4  \ne f'(0)=3$.
   \item $3^x=3$ có nghiệm là $x=1  = a-2$.
   \item  $\log_a9=\log_39=2$.
  \end{itemize}
 } 
\end{ex}

% nguồn: Ok_dungsai_lop11_Lan2
\begin{ex}%[1D7H1-1]%[ma:199d4681]
	Dùng định nghĩa để tính đạo hàm của hàm số $ f(x)=\dfrac{2}{1-x} $ với $ x \neq 1 $. Khi  đó
	\choiceTF 
		{\True Với bất kì $ x_{0} \neq 1 $, ta có $ f'\left(x_{0}\right)=\lim\limits_{x  \rightarrow x_{0}} \dfrac{2}{(1-x)\left(1-x_{0}\right)} $}
		{\True $ f'(2)=2  $}
		 {$ f'(3)=\dfrac{1}{3} $}
		{$ f'(2)+f'(3)=\dfrac{3}{2} $} 
	\loigiai{ 
		\begin{itemchoice}
			\itemch Ta có \begin{eqnarray*}
			f'\left(x_{0}\right) & = &\lim\limits_{x \rightarrow  x_{0}} \dfrac{f(x)-f\left(x_{0}\right)}{x-x_{0}}=\lim\limits_{x \rightarrow x_{0}} \dfrac{\dfrac{2}{1-x}-\dfrac{2}{1-x_{0}}}{x-x_{0}} \\
			& =  &\lim\limits_{x \rightarrow x_{0}} \dfrac{2\left(x-x_{0}\right)}{(1-x)\left(1-x_{0}\right)\left(x-x_{0}\right)}\\
			 &=&\lim\limits_{x \rightarrow x_{0}} \dfrac{2}{(1-x)\left(1-x_{0}\right)}=\dfrac{2}{\left(1-x_{0}\right)^{2}} .
		\end{eqnarray*}
		Vậy $f'(x)=\left(\dfrac{2}{1-x}\right)'=\dfrac{2}{(1-x)^{2}}$ trên các khoảng $(-\infty; 1)$ và  $(1;+\infty)$.
		\itemch  Vì $f'(x) = \dfrac{2}{(1-x)^2}$, $f'(2) = \dfrac{2}{(1-2)^2} = 2$.
		\itemch Vì $f'(x) = \dfrac{2}{(1-x)^2}$, $f'(3) = \dfrac{2}{(1-3)^2} = \dfrac{2}{4} = \dfrac{1}{2}\neq  \dfrac{1}{3}$.
		\itemch Vì  $f'(2) + f'(3) = 2 + \dfrac{1}{2} = \dfrac{5}{2}\neq \dfrac{3}{2}$.
	\end{itemchoice}
} 
\end{ex}

% nguồn: Ok_dungsai_lop11_Lan2
\begin{ex}%[1D7H1-1]%[ma:96527b03]
	Dùng định nghĩa để tính đạo hàm của hàm số $ f(x)=2 x^{3}  $. Khi đó
	 \choiceTF
		{\True Với bất kì $ x_{0}$  ta có $ f'\left(x_{0}\right)=\lim\limits_{x \rightarrow x_{0}} \dfrac{f(x)-f\left(x_{0}\right)}{x-x_{0}} $}
		 {$ f'(1)=-6 $}
		{\True $ f'(0)=0 $} 
		{\True $  f'(2)=24 $}
	\loigiai{ 
		Với  bất kì $ x_{0} $, ta có
		\begin{eqnarray*}
			 f'\left(x_{0}\right) & = & \lim\limits_{x \rightarrow x_{0}} \dfrac{f(x)-f\left(x_{0}\right)}{x-x_{0}}=\lim\limits_{x \rightarrow x_{0}} \dfrac{2 x^{3}-2 x_{0}^{3}}{x-x_{0}} \\
			& = & \lim\limits_{x \rightarrow x_{0}} \dfrac{2\left(x-x_{0}\right)\left(x^{2}+x \cdot  x_{0}+x_{0}^{2}\right)}{x-x_{0}} \\
			& = & \lim\limits_{x \rightarrow x_{0}} 2\left(x^{2}+x \cdot x_{0}+x_{0}^{2}\right) = 6 x_{0}^{2}  .
		\end{eqnarray*}
		Vậy $ f'(x)=\left(2 x^{3}\right)'=6 x^{2} $  trên $ \mathbb{R} $.
		\begin{itemchoice}
			\itemch  Theo định nghĩa.
			\itemch Vì  $f'(1)=6\cdot 1^2=6$.
			\itemch  Vì $f'(0)=6\cdot 0^2=0$.
			\itemch Vì $f'(2)=6\cdot 2^2=24$. 
		\end{itemchoice}
	 }
\end{ex}

% nguồn: Ok_dungsai_lop11_Lan2
\begin{ex}%[1D7H1-1]%[ma:65e2fe91]
	Dùng định nghĩa để tính đạo hàm của hàm số $ f(x)=\dfrac{x-2}{x+1} $ tại điểm $ x_{0}=0 $ ta  được $ f'(0)=a $. Khi đó
	\choiceTF 
		{\True $ f'(0)=\lim\limits_{x \rightarrow 0} \dfrac{f(x)-f(0)}{x-0} $} 
		{$ f'(0)=\lim\limits_{x \rightarrow  0} \dfrac{4}{x+1} $}
		{\True Phương  trình $ 3^{x}=3 $ có nghiệm bằng $ x=a-2 $}
		{$ \log_{a} 9=3 $} 
	 \loigiai{Ta có
		\begin{eqnarray*}
			f'(0) & = & \lim\limits_{x  \rightarrow 0} \dfrac{f(x)-f(0)}{x-0} = \lim\limits_{x \rightarrow 0} \dfrac{\dfrac{x-2}{x+1}-(-2)}{x} \\
			& = & \lim\limits_{x \rightarrow 0} \dfrac{\dfrac{x-2+2(x+1)}{x+1}}{x} = \lim\limits_{x  \rightarrow 0} \dfrac{3x}{x(x+1)} \\
			& = & \lim\limits_{x  \rightarrow 0} \dfrac{3}{x+1} = 3
		\end{eqnarray*}
		Vậy $ f'(0)=3 $. 
		\begin{itemchoice}
			\itemch Theo giải  thích trên.
			\itemch  Vì $\lim\limits_{x \rightarrow 0} \dfrac{3}{x+1}$.
			 \itemch Phương trình $3^x=3$ có nghiệm $x=3-2=1$.
			\itemch Vì  $\log_a9=\log_39=2\neq 3$.
		\end{itemchoice}
	} 
\end{ex}

% nguồn: Ok_dungsai_lop11_Lan2
\begin{ex}%[1D7H1-1]%[ma:c693f935]
	Dùng định  nghĩa tìm đạo hàm của hàm số $ y=x^{2}+1 $ trên tập xác định. Khi đó
	\choiceTF 
		{$ f'\left(x_{0}\right)=\dfrac{f(x)-f\left(x_{0}\right)}{x-x_{0}} $} 
		{\True $  f'\left(x_{0}\right)=\lim\limits_{x \rightarrow x_{0}} \dfrac{f(x)-f\left(x_{0}\right)}{x-x_{0}}=2 x_{0} $}
		{\True $ f'(2)=4 $} 
		{\True  $ f'(2)+f'(-2)=0 $}
	\loigiai{ 
		\begin{itemchoice}
			\itemch Vì $f'(x_0) = \lim\limits_{x \to x_0} \dfrac{f(x)  - f(x_0)}{x - x_0}$.
			\itemch 
			Ta  có
			\begin{eqnarray*}
				f'\left(x_{0}\right)&=&\lim\limits_{x \rightarrow x_{0}}  \dfrac{f(x)-f\left(x_{0}\right)}{x-x_{0}}\\
				 &=&\lim\limits_{x \rightarrow x_{0}} \dfrac{x^{2}+1-\left(x_{0}^{2}+1\right)}{x-x_{0}}\\
				 &=&\lim\limits_{x \rightarrow x_{0}}\left(x+x_{0}\right)=2 x_{0}.
			\end{eqnarray*}
			\itemch Vì $f'(2)=4$; $f'\left(x_{0}\right)=2  x_{0}$ nên $f'(2)=4$.
			\itemch Vì $f'\left(x_{0}\right)=2 x_{0}$  nên $f'(2)+f'\left(-2\right)=4-4=0$.
		\end{itemchoice}
	} 
\end{ex}

% nguồn: Ok_dungsai_lop11_Lan2
\begin{ex}%[1D7H1-1]%[ma:01f7816d]
	Dùng định nghĩa để tính đạo hàm của hàm số $ y=f(x)=x^{2}+2 x $ tại điểm $ x_{0}=1 $. Khi  đó
	 \choiceTF
		{\True $ f'(1)=\lim\limits_{x  \rightarrow 1} \dfrac{f(x)-f(1)}{x-1} $}
		{\True $ f'(1)=\lim\limits_{x \rightarrow 1}  \dfrac{x^{2}+2 x-3}{x-1} $}
		 {$ f'(1)=\lim\limits_{x \rightarrow 1}(x+4) $}
		 {$ f'(1)=a \Rightarrow a>5 $}
	\loigiai{Ta có 
		\begin{eqnarray*}
			f'(1) & = & \lim\limits_{x \rightarrow 1} \dfrac{f(x)-f(1)}{x-1} =  \lim\limits_{x \rightarrow 1} \dfrac{x^{2}+2 x-3}{x-1} \\
			& = &  \lim\limits_{x \rightarrow 1} \dfrac{(x-1)(x+3)}{x-1} = \lim\limits_{x \rightarrow 1}(x+3) = 4
		\end{eqnarray*}
		Vậy $  f'(1)=4 $.
		\begin{itemchoice}
			\itemch Theo giải thích trên. 
			\itemch  Theo giải thích trên.
			\itemch Vì $f'(1)=\lim\limits_{x  \rightarrow 1}(x+3)$.
			\itemch  Vì $f'(1)=4<5$.
		\end{itemchoice}
	} 
\end{ex}

% nguồn: Ok_dungsai_lop11_Lan1
\begin{ex}%[1D7V1-1]%[ma:89b97773]
  Dùng định nghĩa để tính đạo hàm của hàm số $f(x)=2x^3$. Khi đó
 \choiceTF 
   {\True Với bất kì $x_0$ ta có $f'(x_0)=\lim\limits_{x \to x_0}\dfrac{f(x)-f(x_0)}{x-x_0}$}
   {$f'(1)=-6$}
  {\True $f'(0)=0$} 
   {\True $f'(2)=24$}
 \loigiai{ 
  \begin{itemize}
   \item Với bất kì $x_0$  ta có \allowdisplaybreaks
   \begin{eqnarray*}
    f'(x_0) &=&\lim\limits_{x  \to x_0} \dfrac{f(x)-f(x_0)}{x-x_0}\\
    &=&\lim\limits_{x \to x_0}  \dfrac{2x^3-2x_0^3}{x-x_0}\\
    &=&\lim\limits_{x \to x_0} \dfrac{2\left(x-x_0 \right)\left(x^2+x\cdot  x_0+x_0^2 \right)}{x-x_0}\\
    &=&\lim\limits_{x  \to x_0} \dfrac{2\left(x^2+x\cdot x_0+x_0^2 \right)}{1}=6x_0^2.
   \end{eqnarray*}
   Vậy $f'(x)=\left(2x^3 \right)'=6x^2$  trên $\mathbb{R}$.
   \item  $f'(1)=6$.
   \item  $f'(0)=0$.
   \item $f'(2)=24$. 
  \end{itemize}
  }
\end{ex}

% nguồn: Ok_dungsai_lop11_Lan2
\begin{ex}%[1D7V1-1]%[ma:7792ed2f]
	Cho hàm số $f(x)=\heva{&x^2+x&\text{ khi  }x\ge 1\\&ax+b&\text{ khi }x < 1.}$
	 \choiceTF
		{\True $\lim\limits_{x\to  1^+}f(x)=2$}
		{\True  $\lim\limits_{x\to 1^-}f(x)=a+b$}
		{Hàm có đạo hàm tại $x=1$ khi $a=m_0$, $b=n_0$ khi  đó $n_0 > m_0$}
		{\True Hàm có đạo hàm  tại $x=1$ khi $a=m_0$, $b=n_0$ khi đó $n_0+m_0=2$}
	\loigiai{ 
		\begin{itemchoice}
			\itemch Ta  có $\lim\limits_{x\to 1^+}f(x)=\lim\limits_{x\to 1^+}\left(x^2+x\right)=2$.
			\itemch Ta  có $\lim\limits_{x\to 1^-}f(x)=\lim\limits_{x\to 1^-}(ax+b)=a+b$.
			\itemch Hàm có đạo hàm tại $x=1$ thì hàm liên tục  tại $x=1$ $\Leftrightarrow a+b=2$.\hfill$(1)$\\
			 Mặt khác
			\begin{itemize}
				\item $\lim\limits_{x\to 1^+}\dfrac{f(x)-f(1)}{x-1}=\lim\limits_{x\to 1^+}\dfrac{x^2+x-2}{x-1}=\lim\limits_{x\to 1^+}\left(x+2\right)=3$. 
				\item $\lim\limits_{x\to 1^-}\dfrac{f(x)-f(1)}{x-1}=\lim\limits_{x\to 1^-}\dfrac{ax+b-2}{x-1}=\lim\limits_{x\to 1^-}\dfrac{ax-a}{x-1}=a$ (do  $(1)$).
			\end{itemize}
			 Hàm có đạo hàm tại $x=1\Leftrightarrow \lim\limits_{x\to 1^+}\dfrac{f(x)-f(1)}{x-1}=\lim\limits_{x\to 1^-}\dfrac{f(x)-f(1)}{x-1}\Leftrightarrow a=3\Rightarrow b=-1$.
			\itemch Hàm  có đạo hàm tại $x=1$ khi $a=m_0$, $b=n_0$ khi đó $n_0+m_0=2$.
		\end{itemchoice}
	 }
\end{ex}

% nguồn: Ok_dungsai_lop11_Lan2
\begin{ex}%[1D7V1-1]%[ma:3d782076]
	Cho  hàm số $f(x)=\heva{&\dfrac{x^2-1}{x-1}&\text{ khi }x\ne 1\\
	&a&\text{ khi }x=1.}$ 
	 \choiceTF
		{\True Ta  có $\lim\limits_{x\to 1}\dfrac{x^2-1}{x-1}=2$}
		{Với $a=-2$ thì  hàm số có đạo hàm tại $x=1$}
		{\True Với  $a=2$ thì hàm số có đạo hàm tại $x=1$}
		{\True Với $a=m_0$ thì hàm số có đạo  hàm tại $x=1$, khi đó $\lim\limits_{x\to m_0} \left(x^2+2x-3\right)=5$}
	 \loigiai{
		\begin{itemchoice}
			\itemch Ta có $\lim\limits_{x\to 1}\dfrac{x^2-1}{x-1}=\lim\limits_{x\to  1}\dfrac{(x-1)(x+1)}{x-1}=\lim\limits_{x\to 1}(x+1)=2$.
			\itemch Ta  có $f(1)=a$.\\
			Để  hàm số có đạo hàm tại $x=1$ thì $f(x)$ phải liên tục tại $x=1$
			 \[\Leftrightarrow\lim\limits_{x\to 1}f(x)=f(1)\Leftrightarrow a=2.\]
			Vậy với $a=-2$ thì hàm số không liên tục tại $x=1$ nên không có đạo hàm  tại $x=1$.
			\itemch Với $a=2$ thì hàm số liên tục tại $x=1$ nên đạo hàm tại $x=1$  là \[y'(1)=\lim\limits_{x\to 1}\dfrac{f(x)-f(1)}{x-1}=\lim\limits_{x\to 1}\dfrac{x^2-2x+1}{x-1}=\lim\limits_{x\to 1}(x-1)=0.\]
			 Vậy với $a=2$ thì $y'(1)=0$.
			\itemch Với $a=2$ thì  hàm số có đạo hàm tại $x=1$. Khi đó $\lim\limits_{x\to 2}\left(x^2+2x-3\right)=5$.
		\end{itemchoice}
	 }
\end{ex}


%%%%%%%%%%%%%%%% TRẢ LỜI NGẮN %%%%%%%%%%%%%%%%

% nguồn: Ok_TraLoiNgan_Lop11_Lan1
\begin{ex}%[1D7H1-1]%[ma:d6a6e85d]
	Tính đạo hàm của hàm số $f(x)=\sqrt{x^2+1}$ tại $x_0=-1$. (Làm  tròn đến hàng phần mười).
	 \shortans{0,7}
	\loigiai{ 
		Ta có 
		\begin{align*}
			f'(1)&=\lim\limits_{x \rightarrow  1} \dfrac{f(x)-f(1)}{x-1} \\
			&=\lim\limits_{x \rightarrow  1} \dfrac{\sqrt{x^2+1}-\sqrt{2}}{x-1} \\
			 &=\lim\limits_{x \rightarrow 1} \dfrac{x^2-1}{\left(\sqrt{x^2+1}+\sqrt{2}\right)(x-1)} \\
			 &=\lim\limits_{x \rightarrow 1} \dfrac{x+1}{\left(\sqrt{x^2+1}+\sqrt{2}\right)} \\
			 &=\dfrac{1}{\sqrt{2}} \approx 0{,}7.
		\end{align*}
	 }
\end{ex}

% nguồn: Ok_TraLoiNgan_Lop11_Lan1
\begin{ex}%[1D7H1-1]%[ma:53e3b6f2]
	Tính  đạo hàm của hàm số $f(x)=2 x^3+1$ tại $x_0=2$.
	 \shortans{24}
	\loigiai{ 
		Ta  có
		$$ 
		\lim\limits_{x  \rightarrow 2} \dfrac{f(x)-f(2)}{x-2}
		=\lim\limits_{x  \rightarrow 2} \dfrac{2 x^3-16}{x-2}
		 =\lim\limits_{x \rightarrow 2} 2\left(x^2+2 x+4\right)
		 =24.
		$$ 
		Vậy  $f'(2)=24$.
	} 
\end{ex}

% nguồn: Ok_TraLoiNgan_Lop11_Lan2
\begin{ex}%[1D7V1-1]%[ma:11b85b12]
	Tính đạo hàm của hàm số $y=f(x)=\dfrac{1}{x^2+x+1}$ tại $x_0=-2$. (Kết quả làm tròn đến hàng  phần trăm).
	 \shortans{0,33}
	\loigiai{ 
		 Tại $x_0=-2$ ta có
		\begin{eqnarray*}
			\bullet\; f(x)-f\left(x_0  \right)&=&f(x)-f\left(-2\right)=\dfrac{1}{x^2+x+1}-\dfrac{1}{4-2+1}\\
			 &=&\dfrac{1}{3}\cdot\dfrac{3-x^2-x-1}{x^2+x+1}\\
			&=&\dfrac{1}{3}\cdot\dfrac{-x^2-x+2}{x^2+x+1}\\ 
			&=&-\dfrac{1}{3}\cdot\dfrac{\left(x-1\right)\left(x+2\right)}{x^2+x+1}\\ 
			 \bullet\; \dfrac{f(x)-f\left(x_0 \right)}{x-x_0}&=&\dfrac{f(x)-f\left(-2\right)}{x+2}\\
			&=&-\dfrac{1}{3}\cdot\dfrac{\left(x-1\right)\left(x+2\right)}{x^2+x+1}\cdot\dfrac{1}{x+2}\\ 
			&=&-\dfrac{1}{3}\cdot\dfrac{\left(x-1\right)}{x^2+x+1}\\ 
			 \Rightarrow f'(-2)&=&\lim\limits_{x\to x_0}\dfrac{f(x)-f\left(x_0 \right)}{x-x_0}\\
			&=&\lim\limits_{x\to-2}  \left[-\dfrac{1}{3}\cdot\dfrac{\left(x-1\right)}{x^2+x+1} \right]\\
			&=&-\dfrac{1}{3}\cdot\dfrac{\left(-2-1\right)}{\left(-2\right)^2+\left(-2\right)+1}\\ 
			&=&\dfrac{1}{3}  \\
			 &\approx& 0{,}33.
		\end{eqnarray*}
	} 
\end{ex}

% nguồn: Ok_TraLoiNgan_Lop11_Lan2
\begin{ex}%[1D7V1-1]%[ma:68d06b6f]
	Tính đạo hàm của hàm số $f(x)=\sqrt{x^2+1}$ tại $x_0=1$ (kết quả  làm tròn đến hàng phần trăm).
	\shortans{0,71} 
	 \loigiai{
		Ta có \begin{eqnarray*}
		\lim\limits_{x\to  1} \dfrac{f(x)-f(1)}{x-1}&=&\lim\limits_{x\to 1} \dfrac{\sqrt{x^2+1}-\sqrt{2}}{x-1}\\
		&=&\lim\limits_{x\to 1}  \dfrac{x^2-1}{(x-1)\left(\sqrt{x^2+1}+\sqrt{2}\right)}\\
		&=&\lim\limits_{x\to 1}  \dfrac{x+1}{\sqrt{x^2+1}+\sqrt{2}}\\
		&=&\dfrac{\sqrt{2}}{2}. 
	\end{eqnarray*}
	Vậy $f'(1)\approx 0{,}71$. 
 }
\end{ex}

% nguồn: Ok_TraLoiNgan_Lop11_Lan2
\begin{ex}%[1D7V1-1]%[ma:74af7641]
	Hàm số $y=f(x)=x^3-2x$ có đạo  hàm trên tập số thực là $y'=f'(x)=ax^2+bx+c$. Tính $a^2+b-c$.
	\shortans{11} 
	 \loigiai{
		Tại $x_0 \in \mathbb{R}$  tùy ý, ta có
		\begin{itemize}
			\item $f(x)-f\left(x_0  \right)=x^3-2x-x_0 ^3+2x_0=\left(x-x_0 \right)\left(x^2+x\cdot x_0+x_0 ^2-2\right)$.
			\item $\dfrac{f(x)-f\left(x_0 \right)}{x-x_0}=\dfrac{\left(x-x_0 \right)\left(x^2+x\cdot x_0+x_0 ^2-2\right)}{x-x_0}=x^2+x\cdot x_0+x_0  ^2-2$.
			\item $\lim\limits_{x\to x_0}\dfrac{f(x)-f\left(x_0 \right)}{x-x_0}=\lim\limits_{x\to x_0}\left(x^2+x\cdot  x_0+x_0 ^2-2\right)=3x_0 ^2-2$.
		\end{itemize}
		Suy ra $y'=3x^2-2$. Do đó $a=3$,  $b=0$, $c=-2$.\\
		 Vậy $a^2+b-c=11$.
	 }
\end{ex}

% nguồn: Ok_TraLoiNgan_Lop11_Lan2
\begin{ex}%[1D7V1-1]%[ma:38427104]
	Cho hàm số $f(x)=\heva{&2x+1&\text{ khi }x\le 1\\&x^2+bx+1&\text{ khi }x > 1}$. Tìm $b$ để hàm số này có đạo hàm  tại $x=1$.
	 \shortans{1}
	\loigiai{ 
		Ta có 
		\begin{itemize}
			\item  $f(1)=3$;
			 \item $\lim\limits_{x\to 1^-}f(x)=\lim\limits_{x\to 1^-}(2x+1)=3;\lim\limits_{x\to 1^+}f(x)=\lim\limits_{x\to 1^+}\left(x^2+bx+1\right)=b+2$.
		\end{itemize}
		Hàm số $f(x)$ có đạo hàm tại $x=1$ thì $f(x)$ liên tục tại $x=1$. Điều này  tương đương với
		\[\lim\limits_{x\to 1^+}f(x)=\lim\limits_{x\to 1^-}f(x)=f(1)\Rightarrow b+2=3\Rightarrow  b=1.\]
		Vậy với $b=1$ thì $f'\left(1^{+} \right)=f'\left(1^{-}  \right)=3$.
	} 
\end{ex}

% nguồn: Ok_TraLoiNgan_Lop11_Lan2
\begin{ex}%[1D7V1-1]%[ma:a8bdeab2]
	Cho hàm số $f(x)=\heva{&\dfrac{3-\sqrt{4-x}}{4}&\text{  khi }x\ne 0 \\
	&{\dfrac{1}{4}}&\text{ khi }x=0}$. Tính $f'(0)$ (kết quả làm tròn đến hàng phần trăm). 
	 \shortans{0,06}
	 \loigiai{
		Với $x\ne 0$ xét 
		\begin{eqnarray*}
			\lim\limits_{x\to 0}\dfrac{f(x)-f(0)}{x-0}&=&\lim\limits_{x\to 0}\dfrac{\tfrac{3-\sqrt{4-x}}{4}-\tfrac{1}{4}}{x}\\ 
			&=&\lim\limits_{x\to 0}\dfrac{2-\sqrt{4-x}}{4x}\\ 
			 &=&\lim\limits_{x\to 0}\dfrac{1}{4\left(2+\sqrt{4-x}\right)}\\
			 &=&\dfrac{1}{16}.
		\end{eqnarray*}
		Vậy $f'(0)\approx  0{,}06$.
	} 
\end{ex}


%%%%%%%%%%%%%%%% TỰ LUẬN %%%%%%%%%%%%%%%%

% nguồn: Ok_TuanLuan_Toan11_Lan2
\begin{ex}%[1D7N1-1]%[ma:e7467747]
	Tính đạo hàm  của hàm số $y=f(x)=2 x^{3}+x-1 \text { tại } x_{0}=0$. \loigiai{
		 Tại $ x_{0}=0 $ ta có
		\begin{itemize}
			\item $f(x)-f\left(x_{0}\right)=f(x)-f(0)=2  x^{3}+x-1-(-1)=2 x^{3}+x=x\left(2 x^{2}+1\right)$.
			\item $\dfrac{f(x)-f\left(x_{0}\right)}{x-x_{0}}=\dfrac{f(x)-f(0)}{x-0}=\dfrac{x\left(2 x^{2}+1\right)}{x}=2 x^{2}+1$. 
		\end{itemize}
		 Suy ra $f'(0)=\lim\limits_{x \rightarrow x_{0}} \dfrac{f(x)-f\left(x_{0}\right)}{x-x_{0}}=\lim\limits_{x \rightarrow 0}\left(2 x^{2}+1\right)=1$.
	 }
\end{ex}

% nguồn: Ok_TuanLuan_Toan11_Lan2
\begin{ex}%[1D7H1-1]%[ma:f2018dc2]
	Tìm $a$ để hàm số $ f(x)=\heva{&\dfrac{x^{2}-1}{x-1} \text { khi } x \neq 1 \\& a \text { khi  } x=1}$ có đạo hàm tại $ x=1$.
	\loigiai{ 
		Để hàm số có đạo hàm tại $ x=1 $ thì trước hết  $ f(x) $ phải liên tục tại $ x=1 $.\\
		Hay  $ \lim\limits_{x \rightarrow 1} f(x)=\lim\limits_{x \rightarrow 1} \dfrac{x^{2}-1}{x-1}=2=f(1)=a $.\\
		Khi đó, ta có $  \lim\limits_{x \rightarrow 1} \dfrac{f(x)-f(1)}{x-1}=\lim\limits_{x \rightarrow 1} \dfrac{\dfrac{x^{2}-1}{x-1}-2}{x-1}=1 $.\\
		Vậy $  a=2 $ là giá trị cần tìm.
	} 
\end{ex}

% nguồn: Ok_TuanLuan_Toan11_Lan2
\begin{ex}%[1D7H1-1]%[ma:23f6c664]
	Tính đạo hàm của hàm số  sau.
	\begin{enumerate}
		 \item $ y=f(x)=x^{2}-3 x+1 $.
		\item $ y=f(x)=x^{3}-2 x $. 
		\item $ y=f(x)=4 x+3 $. 
	\end{enumerate}
	\loigiai{ 
		\begin{enumerate}
			 \item Tại $ x_{0} \in \mathbb{R} $ tùy ý, ta có
			\begin{align*}
				f(x)-f\left(x_{0}\right)&=x^{2}-3  x+1-x_{0}^{2}+3 x_{0}-1\\
				 &=\left(x-x_{0}\right)\left(x+x_{0}-3\right) \\
				\dfrac{f(x)-f\left(x_{0}\right)}{x-x_{0}}&=\dfrac{\left(x-x_{0}\right)\left(x+x_{0}-3\right)}{x-x_{0}}=x+x_{0}-3 \\ 
				\lim\limits_{x \rightarrow x_{0}}  \dfrac{f(x)-f\left(x_{0}\right)}{x-x_{0}}&=\lim\limits_{x \rightarrow x_{0}}\left(x+x_{0}-3\right)=2 x_{0}-3.
			\end{align*}
			Suy ra  $y'=2 x-3$.
			\item Tại $ x_{0} \in \mathbb{R}  $ tùy ý, ta có
			\begin{align*}
				 f(x)-f\left(x_{0}\right)&=x^{3}-2 x-x_{0}^{3}+2 x_{0}\\
				&=\left(x-x_{0}\right)\left(x^{2}+x \cdot x_{0}+x_{0}^{2}-2\right) \\ 
				\dfrac{f(x)-f\left(x_{0}\right)}{x-x_{0}}&=\dfrac{\left(x-x_{0}\right)\left(x^{2}+x \cdot  x_{0}+x_{0}^{2}-2\right)}{x-x_{0}}=\left(x^{2}+x \cdot x_{0}+x_{0}^{2}-2\right) \\
				 \lim\limits_{x \rightarrow x_{0}} \dfrac{f(x)-f\left(x_{0}\right)}{x-x_{0}}&=\lim\limits_{x \rightarrow x_{0}}\left(x^{2}+x \cdot x_{0}+x_{0}^{2}-2\right)=3 x_{0}^{2}-2.
			\end{align*}
			Suy  ra $y'=3 x^{2}-2$.
			\item Tại $ x_{0} \in \mathbb{R} $ tùy ý, ta có 
			\begin{align*}
				f(x)-f\left(x_{0}\right)&=4 x+3-4 x_{0}+3=4\left(x-x_{0}\right)  \\
				 \dfrac{f(x)-f\left(x_{0}\right)}{x-x_{0}}&=\dfrac{4\left(x-x_{0}\right)}{x-x_{0}}=4 \\
				\lim\limits_{x \rightarrow x_{0}} \dfrac{f(x)-f\left(x_{0}\right)}{x-x_{0}}&=\lim\limits_{x \rightarrow x_{0}}  4=4.
			\end{align*}
			Suy ra $y'=4$. 
		\end{enumerate}
	 }
\end{ex}

% nguồn: Ok_TuanLuan_Toan11_Lan2
\begin{ex}%[1D7H1-1]%[ma:47457d23]
	 Tính đạo hàm của các hàm số sau tại các điểm đã chỉ ra.
	\begin{enumerate}
		\item $ f(x)=2 x^{3}+1 $  tại $ x=2$.
		\item $ f(x)=\sqrt{x^{2}+1}  $ tại $ x=1$.
		\item $f(x)=\heva{&\dfrac{\sqrt{x^{3}+x^{2}+1}-1}{x} \text { khi } x \neq 0 \\& 0  \quad \text { khi } x=0.}\text { tại } x=0 $.
	\end{enumerate}
	 \loigiai{
		\begin{enumerate}
			\item Ta có $\lim\limits_{x \rightarrow 2} \dfrac{f(x)-f(2)}{x-2}=\lim\limits_{x \rightarrow  2} \dfrac{2 x^{3}-16}{x-2}=\lim\limits_{x \rightarrow 2} 2\left(x^{2}+2 x+4\right)=24$.\\
			Suy ra  $f'(2)=24$.
			 \item Ta có
			\begin{eqnarray*}
				\lim\limits_{x \rightarrow 1} \dfrac{f(x)-f(1)}{x-1}&=&\lim\limits_{x \rightarrow 1} \dfrac{\sqrt{x^{2}+1}-\sqrt{2}}{x-1} \\ 
				 &=&\lim\limits_{x \rightarrow 1} \dfrac{(x-1)(x+1)}{(x-1)\left(\sqrt{x^{2}+1}+\sqrt{2}\right)}=\dfrac{1}{\sqrt{2}}.
			\end{eqnarray*}
			Suy ra $f'(1)=\dfrac{1}{\sqrt{2}}$. 
			 \item Ta có $ f(0)=0 $, do đó
			\begin{eqnarray*}
				\lim\limits_{x  \rightarrow 0} \dfrac{f(x)-f(0)}{x}&=&\lim\limits_{x \rightarrow 0} \dfrac{\sqrt{x^{3}+x^{2}+1}-1}{x^{2}}\\
				 &=&\lim\limits_{x \rightarrow 0} \dfrac{x+1}{\sqrt{x^{3}+x^{2}+1}+1}=\dfrac{1}{2}.
			\end{eqnarray*}
			 Vậy $f'(0)=\dfrac{1}{2}$.
		\end{enumerate}
	 }
\end{ex}

% nguồn: Ok_TuanLuan_Toan11_Lan2
\begin{ex}%[1D7H1-1]%[ma:cc02a08c]
	Tính đạo  hàm của các hàm số tại $x_0$ dưới đây.
	\begin{enumerate}
		\item $ y=f(x)=\dfrac{1}{x^{2}+x+1} $ tại $ x_{0}=-2  $.
		\item $  y=f(x)=\dfrac{x^{2}+x-3}{2 x-1} $ tại $ x_{0}=3 $.
	\end{enumerate}
	\loigiai{ 
		\begin{enumerate}
			 \item Tại $ x_{0}=-2 $ ta có
			\begin{eqnarray*}
				f(x)-f\left(x_{0}\right)&=&f(x)-f(-2)=\dfrac{1}{x^{2}+x+1}-\dfrac{1}{4-2+1} \\ 
				&=&\dfrac{1}{3} \cdot \dfrac{-x^{2}-x+2}{x^{2}+x+1}=-\dfrac{1}{3} \cdot  \dfrac{(x-1)(x+2)}{x^{2}+x+1} \\
				\dfrac{f(x)-f\left(x_{0}\right)}{x-x_{0}}&=&\dfrac{f(x)-f(-2)}{x+2}=-\dfrac{1}{3} \cdot \dfrac{(x-1)}{x^{2}+x+1}. 
			\end{eqnarray*}
			 Suy ra
			\[f'(-2)=\lim\limits_{x \rightarrow x_{0}}  \dfrac{f(x)-f\left(x_{0}\right)}{x-x_{0}}=\lim\limits_{x \rightarrow-2}\left[-\dfrac{1}{3} \cdot \dfrac{(x-1)}{x^{2}+x+1}\right]=\dfrac{1}{3}.\]
			\item Tại $ x_{0}=3 $  ta có
			\begin{eqnarray*}
				f(x)-f(3)&=&\dfrac{x^{2}+x-3}{2 x-1}-\dfrac{9}{5}=\dfrac{5 x^{2}-13 x-6}{5(2 x-1)}=\dfrac{(x-3)(5  x+2)}{5(2 x-1)} \\
				\dfrac{f(x)-f\left(x_{0}\right)}{x-x_{0}}&=&\dfrac{f(x)-f(3)}{x-3}=\dfrac{(x-3)(5 x+2)}{5(2 x-1)} \cdot  \dfrac{1}{x-3}=\dfrac{(5 x+2)}{5(2 x-1)} \\
				\lim\limits_{x \rightarrow x_{0}} \dfrac{f(x)-f\left(x_{0}\right)}{x-x_{0}}&=&\lim\limits_{x \rightarrow  3} \dfrac{f(x)-f(3)}{x-3}=\lim\limits_{x \rightarrow 3} \dfrac{(5 x+2)}{5(2 x-1)}=\dfrac{17}{25}.
			\end{eqnarray*}
			 Suy ra $f'(3)=\dfrac{17}{25}$.
		\end{enumerate}
	 }
\end{ex}
